\chapter{Computation, data and the statistical protocol}\label{ch:computation}
\chaptermark{Computation and data}

The paper's theorems are proved; its conjectures, its sharpness claims and a dozen of its smaller statements rest on computations with PARI/GP on
Cremona's tables, organised by a protocol of working sets and hold-outs. This chapter makes that side of the paper as checkable as the other. It
teaches PARI/GP and the \texttt{ecdata} files from zero, with exactly the functions the paper uses; reads the paper's appendix code line by line and
proves what it computes; runs it on the running examples; explains the protocol and judges the statistics, including the calibration argument
of Section sec:sharp; and reproduces every table of the paper on all optimal curves of conductor below $10{,}000$, from the very data commit the paper
names, with scripts that run in minutes. Along the way it records the places where the paper's data sentences need correcting.

\section*{Core results at a glance}

Throughout, $E$ is an optimal curve given by Cremona's minimal model, $\ell$ an odd prime, $q$ a bad prime, $v=v_\ell(\degphi)$; outputs are those of PARI/GP
2.17.2 through \texttt{cypari2} 2.2.4, and ``below $10{,}000$'' means the $38{,}042$ optimal curves of conductor at most $9{,}999$ in the files of \texttt{ecdata}
commit \texttt{25cec5e}. Items tagged \emph{(proved here, by computation)} are established by the computation printed or described in the item, and
rely on PARI computing what its documentation says.

\begin{enumerate}[label=\textbf{\thechapter.\arabic*}]
\item \textbf{Versions and data.} PARI/GP 2.17.2 via \texttt{cypari2} 2.2.4 (the paper: 2.15, and 2.17.2 for its checks below $10{,}000$); the six files
\texttt{allcurves}, \texttt{allbsd}, \texttt{allisog}, \texttt{alllabels}, \texttt{degphi}, \texttt{galrep} for conductors $0$ to $9{,}999$, identical (SHA-256) with commit
\texttt{25cec5e} of \texttt{ecdata}, the paper's. \emph{(convention)} \Cref{conv:computation:versions,conv:computation:checklist}.
\item \textbf{The five bounds.} $B^{\mathrm{nc}}_\ell=\sum_{q\ne\ell}e_\ell(E,q)+\sum_qt_\ell(E,q)+[\ell=3]w_3(E)$; $B^{\mathrm{rule}}_\ell=B^{\mathrm{nc}}_\ell-h_\ell(E)$ (the right side of (eq:rule));
$B^{\mathrm{III}}_\ell$ is $B^{\mathrm{nc}}_\ell$ without the $e_\ell$ of the types $\kod{III}$, $\kod{III}^*$ at $q\ne\ell$; $B^{\mathrm{main}}_\ell=\sum_qv_\ell(\#\Phig_q)$; $B^{\mathrm{tw}}_\ell=\sum_qw_\ell(E,q)$.
Violation: $v<b$; equality: $v=b>0$; excess $v-b$; rates among $b>0$. \emph{(convention)} \Cref{conv:computation:bounds}.
\item \textbf{Populations.} CM by the thirteen CM $j$-invariants; $E[\ell]$ reducible iff a \texttt{galrep} label begins with $\ell$B or $\ell$Cs; base population at $\ell$:
non-CM, $E[\ell]$ irreducible; Eisenstein: non-CM, reducible; hypothesis population: base, $\ell\nmid N$, at $\ell=3$ image not $\Ns$. Below $10{,}000$: $38{,}042$ optimal
curves, $189$ CM; base populations $33{,}851$, $37{,}319$, $37{,}691$, $37{,}839$, $37{,}823$ at $\ell=3,5,7,11,13$. \emph{(convention; counts proved here, by computation)}
\Cref{conv:computation:populations}.
\item \textbf{Running PARI; the language; pitfalls.} \texttt{gp} or \texttt{cypari2}; without \texttt{elldata}, \texttt{ellsearch} fails, so files are read directly;
default stack $8{,}000{,}000$ bytes, maximum $8{,}003{,}584$; \texttt{factor} drops the content; a GP file must be loaded with \texttt{read};
\texttt{mfsturm} of an \texttt{mfinit} space is smaller than Sturm's bound $k\mu(N)/12$ ($50$ against $112$ at level $539$). \emph{(stated; reference
\cite{PARIusers}; behaviour checked here)} \Cref{rem:computation:install,rem:computation:gp-language,ex:computation:first-session,warn:computation:pitfalls}.
\item \textbf{\texttt{ellglobalred}.} Returns $[N,[u,r,s,t],\prod c_p,F,L]$ with $F$ the factorisation of $N$ (so \texttt{R[4][,1]} is the column of bad primes) and
$L[i]=[f_p,k_p,0,c_p]$, $k_p$ PARI's Kodaira code; for \curve{726a1}, $L=[[1,10,0,2],[1,7,0,1],[2,3,0,2]]$ (types $\kod{I}_6$, $\kod{I}_3$, $\kod{III}$).
\emph{(definition)} \Cref{def:computation:ellglobalred}.
\item \textbf{Kodaira codes.} $1$: $\kod{I}_0$; $n+4$: $\kod{I}_n$ and $-(n+4)$: $\kod{I}_n^*$, for $n\ge1$ only; $-1$: $\kod{I}_0^*$; $\pm2,\pm3,\pm4$: $\kod{II},\kod{III},\kod{IV}$ and their duals.
\texttt{k >= 5} iff multiplicative; \texttt{abs(k) == 3} iff $\kod{III}$ or $\kod{III}^*$. \emph{(convention)} \Cref{conv:computation:code-idioms}.
\item \textbf{The symmetric square.} \texttt{lfuneuler(lfunsympow(E, 2), q)} $=1/P_q(X)$, so $P_q(q^{-2})=L_q(A_f,1)^{-1}$ is \texttt{subst(1/lfuneuler(S, q), x, 1/q\^{}2)};
for \curve{121a1}, \texttt{lfunparams(S)} $=$ \texttt{[121, 3, [0, 0, 1]]} and $P_{11}=1+11X$, $P_{11}(11^{-2})=12/11$. \emph{(proved here, by computation)}
\Cref{rem:computation:lfunsympow}.
\item \textbf{Lemma lem:euler against PARI.} The factor read off the Kodaira type equals PARI's $P_q(q^{-2})$ at all $67{,}745$ pairs (optimal curve below $10{,}000$, bad
$q\ge5$). \emph{(proved here, by computation)} \Cref{ex:computation:euler-by-hand}.
\item \textbf{Tate parameters.} \texttt{ellinit(E, O(q\^{}k)).tate[3]} is $q_E$ (not \texttt{ellinit(E, q)}, which reduces modulo $q$, and not \texttt{elltatepairing});
$\bar u=10$ for \curve{11a1} at $11$, $5$ for \curve{77c1} and \curve{539c1} at $7$, $6$ for \curve{791a1} at $7$, $8$ for \curve{1577a1} at $19$; wrong at $q=2$ when $v_2(j)=-1$
(PARI 2.17.2; checked on \curve{208a1}). \emph{(proved here, by computation)} \Cref{ex:computation:tate-parameter}.
\item \textbf{\texttt{ns3} and \texttt{galrep}.} For \curve{726a1}, $\psi_3=(x^2-x-74)(3x^2+8x+20)$. Among the $37{,}853$ non-CM optimal curves below $10{,}000$, \texttt{ns3} $=1$ exactly on the
$90$ labelled \texttt{3Ns}; among the CM ones on $13$ curves with $j\in\{-32768,8000\}$, unlabelled in \texttt{galrep}. \emph{(proved here, by computation)}
\Cref{ex:computation:mod3-image}.
\item \textbf{Lemma lem:twist against PARI.} All $85{,}848$ pairs (optimal curve below $10{,}000$, odd bad $q$) agree with Lemma lem:twist(i), (ii); the exchanges at $q=3$
hold in all $5{,}122$ pairs. \emph{(proved here, by computation)} \Cref{ex:computation:twists}.
\item \textbf{\texttt{ellmoddegree} and \texttt{ellweilcurve}.} \texttt{ellmoddegree(E)} $=\deg\varphi/c_\varphi^2$ ($1$, $5$, $1/5$ on \curve{11a1}, \curve{11a2}, \curve{11a3}); \texttt{ellweilcurve}
marks the optimal curve by the Smith invariants $[1,1]$ (\curve{11a1}, \curve{990h3}); the largest degrees below $10{,}000$ need a stack of $16$ to $32$ MB. \emph{(stated;
reference \cite{PARIusers}; examples proved here, by computation)} \Cref{rem:computation:ellmoddegree}.
\item \textbf{The \texttt{ecdata} files.} Formats of \texttt{allcurves}, \texttt{allbsd}, \texttt{allisog}, \texttt{alllabels}, \texttt{degphi}, \texttt{galrep}; these six do not contain
per-prime Kodaira symbols or Tamagawa numbers, the repository's \texttt{curvedata} does. \emph{(definition)} \Cref{def:computation:ecdata-files,ex:computation:reading-ecdata}.
\item \textbf{Cross-checks with \texttt{curvedata}, \texttt{opt\_man}, \texttt{alldegphi}.} Kodaira codes, $c_p$, $f_p$ agree with \texttt{ellglobalred} on all $64{,}687$
curves; exactly one optimal curve per class, the one of \texttt{degphi}, with Manin constant $1$. \emph{(proved here, by computation)} \Cref{ex:computation:curvedata}.
\item \textbf{Optimality.} Curve number $1$ is optimal in every class below $10{,}000$ except \curve{990h}, whose optimal curve is \curve{990h3}; for $E[\ell]$ irreducible,
$v_\ell$ of the degree does not depend on the curve of the class. \emph{(proved here, by computation; the second statement stated, reference \Cref{ch:moddeg})}
\Cref{rem:computation:optimality-data}.
\item \textbf{Labels.} Cremona and LMFDB labels (\curve{11a1} $=$ \texttt{11.a2}); Sutherland's labels, reducible iff $\ell$B or $\ell$Cs; counts of the $3$-adic labels
among non-CM optimal curves below $10{,}000$. \emph{(convention)} \Cref{conv:computation:labels}.
\item \textbf{The appendix code.} The functions \texttt{rule} and \texttt{ns3}, verbatim, with a line-by-line reading. \emph{(construction)}
\Cref{con:computation:rule-listing,rem:computation:rule-reading}.
\item \textbf{What \texttt{rule} computes.} Given PARI's documented behaviour: \texttt{e} $=e_\ell(E,q)$ for $q\ne\ell$, \texttt{t} $=t_\ell(E,q)$ for all $q$, \texttt{e} $=w_3(E)$ at $q=\ell=3$;
\texttt{ns3(E)} $=1$ iff $\psi_3$ is a product of two irreducible quadratics, and \texttt{ns3} $=h_3$ when $E[3]$ is irreducible; for the optimal curve (and $E[3]$
irreducible if $\ell=3$), \texttt{rule(E, l, 1)} $=[v,B^{\mathrm{rule}}_\ell]$ and \texttt{rule(E, l)} $=[v,B^{\mathrm{III}}_\ell]$. \emph{(proved here, given PARI's documented behaviour and
\Cref{ch:galrep})} \Cref{prop:computation:rule-correct}.
\item \textbf{Where \texttt{rule} misleads.} Non-optimal models (\curve{11a3}: \texttt{[-1, 0]}; \curve{990h1}: \texttt{[4, 3]}), CM curves, Eisenstein curves (\curve{11a1} at $5$:
\texttt{[0, 1]}), $\ell\mid N$, codes at $n=0$, versions. \emph{(proved here, by computation)} \Cref{warn:computation:rule-pitfalls}.
\item \textbf{Sessions.} \texttt{rule} returns \texttt{[1, 1]} on \curve{121a1}; \texttt{[2, 2]} on \curve{726a1} with and without \texttt{iii} ($B^{\mathrm{nc}}_3=3$, \texttt{ns3} $=1$); \texttt{[2, 2]} on
\curve{539c1}; \texttt{[2, 2]} and \texttt{[1, 1]} on \curve{980i1} at $3$, $5$; \texttt{[0, 0]} on \curve{99a1} ($\degphi=4$, $r_E=12$); the prime $2$ contributes through PARI's factor on
\curve{336f1}, \curve{1232l1} ($\ell=3$) and \curve{592c1} ($\ell=5$); \curve{11a1} at $5$ gives \texttt{[0, 1]}, outside every statement; the Tamagawa refinement is attained on
\curve{791a1} and \curve{1577a1}. \emph{(proved here, by computation)}
\Cref{ex:computation:121a1,ex:computation:726a1,ex:computation:539c1,ex:computation:980i1,ex:computation:99a1,ex:computation:wild-two,ex:computation:11a1-eisenstein,ex:computation:tamagawa-session}.
\item \textbf{Modular degrees.} Zagier's, Cremona's and Watkins's methods; \texttt{ellmoddegree} implements Watkins's; it agrees with \texttt{degphi} on all $38{,}042$ optimal
curves below $10{,}000$. PARI does not compute $r_E$. \emph{(methods stated, reference \cite{Watkins02,Cremona97}; the check proved here, by computation)}
\Cref{rem:computation:moddeg-methods,ex:computation:moddeg-check,rem:computation:congruence-number}.
\item \textbf{Units and rates.} One optimal curve per class at a fixed $\ell$; equality rate and mean excess among positive bounds; base rates. \emph{(definition)}
\Cref{def:computation:units,def:computation:rates}.
\item \textbf{Base rates and pooling.} At $\ell=3$ below $10{,}000$, $3\mid\degphi$ for $76.1\%$ of the base population, for $100\%$ of the curves with a prime $q\ge5$ of type
$\kod{II}$ to $\kod{IV}^*$, $\kod{I}_n^*$ or $\kod{I}_n$ with $3\mid n$, and for $73\%$ to $76\%$ for the other types; pooled rates mix ranges. \emph{(proved here, by computation)}
\Cref{ex:computation:base-rates,rem:computation:simpson}.
\item \textbf{The protocol.} The four ranges; the status of each statement: (eq:main), Conjecture conj:eis blind-tested on the hold-out; (eq:twist) blind-tested if
fixed before the hold-out counts of (eq:main) were seen; the Euler terms blind-tested only on the out-of-table curves; the weights $w_3$, the correction
$h_3$ and the III-free form never blind-tested. \emph{(a reading of the paper's text)}
\Cref{rem:computation:protocol,rem:computation:test-status,rem:computation:what-is-tested,pl:computation:protocol}.
\item \textbf{Multiple comparisons.} $m$ false candidates surviving with probabilities $\pi$, $\pi'$ leave about $m\pi$ and $m\pi\pi'$; the risk sits in the choice of
form; features read off the data are validated only by other data. \emph{(proved here)} \Cref{rem:computation:multiple-comparisons}.
\item \textbf{The protocol in miniature.} The minima of $v_3(\degphi)$ at $q=3$ on $N<5000$ are the paper's $w_3$ and hold on $5000\le N<10000$ ($11{,}699$ curves with $3\mid N$,
no violation). \emph{(proved here, by computation)} \Cref{ex:computation:toy-protocol}.
\item \textbf{The null model.} Under $H_0$ ($v$ independent of the bound, common law $F$): $E[X]=\sum_{b>0}F(v<b)$ without independence; with independence
$\operatorname{Var}X\le E[X]$ and $P(X=0)\le e^{-E[X]}$; with pairwise independence $P(X=0)\le1/E[X]$; the pooled and the bound-$0$ laws of $v$ agree under $H_0$.
\emph{(proved here)} \Cref{prop:computation:null-model}.
\item \textbf{Calibration.} The paper: $605{,}000$ expected violations, $0$ observed. Below $10{,}000$ at $\ell=3$: $\hat E=14{,}275$ (standard deviation $61.5$), stratified
$14{,}177$, permutation mean $7{,}638$; $0$ observed; $v=0$ for $59.7\%$ of bound-$0$ curves and $23.9\%$ of all. It shows strength, not validity.
\emph{(the paper's numbers stated, reference Section sec:sharp; the book's proved here, by computation)}
\Cref{ex:computation:calibration-paper,ex:computation:calibration-small,rem:computation:calibration-reading}.
\item \textbf{Sharpness.} Below $10{,}000$: $14{,}420$ equalities at $\ell=3$ against $4{,}548$ expected under $H_0$; equality rates $71.0\%$, $84.5\%$, $90.7\%$, $96.5\%$, $96.0\%$ at
$\ell=3,\dots,13$; the odd part of $\degphi$ equals the product of the twist bounds for $17.1\%$ of $33{,}137$ curves, and the quotient is always an integer.
\emph{(proved here, by computation)} \Cref{ex:computation:sharpness,ex:computation:odd-part}.
\item \textbf{The tables in small.} Below $10{,}000$: no violation of (eq:rule), of the III-free form, of (eq:main), of (eq:twist) or of Conjecture conj:eis at
$\ell=3,5,7,11,13$; the $36$ failures without the correction are $\Ns$ curves, by exactly one ($24$ if (eq:thm) is read literally); the minima of Section sec:iso
equal $e_\ell$ cell by cell; the Tamagawa refinement has no exception ($1{,}050$ and $252$ curves at $\ell=3,5$, as in the paper); strict twist inequalities in
cases (a) and (b). \emph{(proved here, by computation)}
\Cref{ex:computation:repro-uniform,ex:computation:repro-counts,ex:computation:repro-ns,ex:computation:repro-iso,ex:computation:repro-tam,ex:computation:repro-eis,ex:computation:repro-twist,ex:computation:repro-further,rem:computation:comparison}.
\item \textbf{A twist relation at $q=\ell$.} Below $10{,}000$, for $\ell\le13$: a curve of type $\kod{II}^*$, $\kod{III}^*$ or $\kod{IV}^*$ at $\ell$ has $v_\ell$ of its degree exactly one more than its
twist by $\chi_{\ell^*}$ ($4{,}412$ cases); type $\kod{I}_n^*$ gives $0$; type $\kod{I}_0^*$ gives $v_\ell((\ell+1)^2-a_\ell^2)$ when $E[\ell]$ is irreducible. \emph{(an observation, by
computation; no proof)} \Cref{ex:computation:twist-at-ell}.
\item \textbf{Reproducibility.} The paper's repository (described, not used); the book's scripts \texttt{sessions.gp}, \texttt{table\_small.gp} and the Python
scripts; a checklist: pin the data, record versions, separate data and results, rerun, cross-check, state what was not checked. \emph{(construction;
convention)} \Cref{rem:computation:paper-repository,con:computation:sessions,con:computation:table-small,rem:computation:book-scripts,conv:computation:checklist}.
\end{enumerate}

\section*{Where the paper uses this chapter}
\begin{sloppypar}

\begin{itemize}
\item Section 1.1 (``PARI's \texttt{lfunsympow(E,2)} returns these factors at every prime, including $2$ and $3$''): \Cref{rem:computation:lfunsympow} says what the
function returns and how $P_q(q^{-2})$ and $e_\ell(E,q)$ are read off; \Cref{ex:computation:wild-two} runs it at $q=2$.
\item Definition def:terms and Conjecture conj:euler: the five bounds of \Cref{conv:computation:bounds}; the term $t_\ell$ at $q=\ell$ as the code computes it
(\Cref{rem:computation:rule-reading}(g)).
\item Section 1.4 (component groups; optimality): \Cref{def:computation:ellglobalred,conv:computation:code-idioms,rem:computation:optimality-data}.
\item Lemma lem:euler (``In each case the values were confirmed against PARI's \texttt{lfunsympow} on every curve of the tables''): \Cref{ex:computation:euler-by-hand}.
\item Lemma lem:twist and the appendix's \texttt{twist\_lemma\_check.gp}: \Cref{ex:computation:twists}.
\item Remark rem:wild and Proposition prop:shapes at $q=2,3$: \Cref{ex:computation:wild-two}, with the population of the counts corrected.
\item Proposition prop:ns (the $3$-division polynomial) and the function \texttt{ns3}: \Cref{ex:computation:mod3-image}, \Cref{prop:computation:rule-correct}(ii).
\item Proposition prop:forms (``PARI's factor at $2$ is $1+2X$ for every one of them''): \Cref{ex:computation:repro-further}(c).
\item Lemma lem:tamfactor and Section sec:tam (the Tate parameter): \Cref{ex:computation:tate-parameter,ex:computation:tamagawa-session,ex:computation:repro-tam}.
\item The Example of Section sec:proof and the Remark after it (99A1, 726a1, 539d1): \Cref{ex:computation:121a1,ex:computation:539c1,ex:computation:980i1,ex:computation:99a1,rem:computation:congruence-number}.
\item Section sec:data and all its subsections and tables: \Cref{sec:computation:protocol,sec:computation:calibration,sec:computation:repro}; the status of each data statement is
\Cref{rem:computation:test-status}; the calibration argument of Section sec:sharp is \Cref{prop:computation:null-model}.
\item The paragraph after Theorem thm:twist (the twist pairs): \Cref{ex:computation:repro-twist}.
\item Section sec:two (the $2$-adic inequality; Calegari-Emerton and CM curves): \Cref{ex:computation:repro-further}(d), (e).
\item The appendix (the functions \texttt{rule} and \texttt{ns3}, the scripts, the memory remark): \Cref{sec:computation:rule,sec:computation:sessions,sec:computation:reproducibility}.
\end{itemize}

The mathematics behind every computation is in earlier chapters: Tate's algorithm and the Kodaira codes in \Cref{ch:reduction}, the mod $3$ images in
\Cref{ch:galrep}, optimality, Manin constants and the congruence number in \Cref{ch:moddeg}, the symmetric square and its local factors in
\Cref{ch:lfunctions}, the Tamagawa factor of the adjoint in \Cref{ch:blochkato}, what each data statement of the paper means in \Cref{ch:paper},
and the vocabulary of elementary probability and statistics in \Cref{ch:prereq}. This chapter computes, counts and judges; it cites the
mathematics and does not repeat it.

\section{Introduction and conventions}\label{sec:computation:intro}

The paper rests on computations as much as on proofs. Its Section sec:data reports tests of its inequalities on every optimal curve of conductor up to $500000$ whose
optimality the tables determine, and on $893$ curves beyond; its appendix prints a PARI/GP function, \texttt{rule(E, l, iii)}, that checks one curve; a dozen sentences elsewhere appeal to
PARI (``PARI's \texttt{lfunsympow(E,2)} returns these factors at every prime'', ``the values were confirmed against PARI's
\texttt{lfunsympow} on every curve of the tables''). This chapter has four jobs. It teaches the software and the data from zero: PARI/GP
with exactly the functions the paper uses, and Cremona's tables in the form of the \texttt{ecdata} files (Sections
\ref{sec:computation:pari} to \ref{sec:computation:ecdata}). It reads the appendix code line by line, proves that it computes what the paper
says under stated conditions, and runs it on the running examples, including the cases where it gives an answer that is not the paper's
quantity (Sections \ref{sec:computation:rule} to \ref{sec:computation:moddeg}). It explains the paper's protocol and judges its statistics:
which statement was blind-tested and which only confirmed, what the calibration argument shows and what it does not
(Sections \ref{sec:computation:populations} to \ref{sec:computation:calibration}). And it reproduces every table of the paper in small, on all
$38{,}042$ optimal curves of conductor below $10{,}000$, with scripts the reader can rerun in minutes (Sections \ref{sec:computation:repro}
and \ref{sec:computation:reproducibility}).

Two kinds of numbers appear. Numbers the book computed are printed with the range they cover and were produced by the code printed or
described here; numbers the paper reports for its large ranges are quoted as the paper's and marked as such: they cannot be rerun from the
data used in this book. The book's ranges were not blind (every statement was known before the book's data were looked at), so the book's
counts are reproductions, not tests.

\begin{convention}[software, data and outputs]\label{conv:computation:versions}
\begin{enumerate}[label=(\alph*)]
\item \emph{Software.} The paper used PARI/GP 2.15 \cite{PARI} on its large ranges and PARI/GP 2.17.2 \cite{PARI217} for the checks of its
revised version below conductor $10{,}000$ (Section sec:data). Every output printed in this chapter was produced with PARI/GP 2.17.2
(\texttt{version()} returns \texttt{[2, 17, 2]}) through the Python interface \texttt{cypari2} 2.2.4 \cite{cypari2}, under Python 3.13 on
Linux; no stand-alone \texttt{gp} binary was used. Every GP function printed here was loaded with \texttt{read} and run as printed.
\item \emph{Data.} Six files of the \texttt{ecdata} repository \cite{ecdata} for conductors $0$ to $9{,}999$:
\texttt{allcurves}, \texttt{allbsd}, \texttt{allisog}, \texttt{alllabels}, \texttt{degphi}, \texttt{galrep} (each with the suffix
\texttt{.00000-09999}). Their SHA-256 sums are listed in \Cref{conv:computation:checklist}; they agree with the files of commit
\texttt{25cec5e} of the repository (full hash \texttt{25cec5ecfec8b9f016eb1631ac633194c2bed39f}, dated 4 February 2026), the commit the
paper names. Four further files of the same commit and range (\texttt{alldegphi}, \texttt{curvedata}, \texttt{classdata},
\texttt{opt\_man}) were used for the cross-checks of \Cref{ex:computation:curvedata}.
\item \emph{Outputs.} A listing that shows output shows it exactly as PARI 2.17.2 prints it; inputs are preceded by \texttt{?} as in a
\texttt{gp} session. Where PARI 2.15 might behave differently we say so (\Cref{warn:computation:rule-pitfalls}(f)); we have not run 2.15.
\item \emph{Models and labels.} Curves are given by Cremona's reduced global minimal models $[a_1,a_2,a_3,a_4,a_6]$ and named by Cremona's
labels, \curve{121a1}; ``optimal'' means $X_0(N)$-optimal (\Cref{ch:moddeg}). The LMFDB's labels appear only in
\Cref{conv:computation:labels}.
\item \emph{Letters.} As in the paper, $q$ is a bad prime and $\ell$ an odd prime; $q_E$ is a Tate parameter. In code, \texttt{x} is
PARI's polynomial variable, which in the output of \texttt{lfuneuler} stands for $q^{-s}$; we write $X$ for it in the text, as the paper
does in $P_q(X)$.
\end{enumerate}
\end{convention}

\begin{convention}[the five bounds; violation, equality, excess]\label{conv:computation:bounds}
Let $E$ be an optimal curve, $\ell$ an odd prime, and $v=v_\ell(\degphi)$ the \emph{observed valuation}. With $e_\ell(E,q)$, $t_\ell(E,q)$
($q=\ell$ included), $w_3(E)$ and $h_\ell(E)$ as in Definition def:terms and Conjecture conj:euler (\Cref{ch:lfunctions,ch:paper}), with $\#\Phig_q$
the order of the geometric component group at $q$ and $w_\ell(E,q)$ the twist-symmetrised weight of the paper's Section 1.4 (\Cref{ch:reduction}),
put
\begin{align*}
B^{\mathrm{nc}}_\ell(E)&=\sum_{q\mid N,\,q\ne\ell}e_\ell(E,q)+\sum_{q\mid N}t_\ell(E,q)+[\ell=3]\,w_3(E), &
B^{\mathrm{rule}}_\ell(E)&=B^{\mathrm{nc}}_\ell(E)-h_\ell(E),\\
B^{\mathrm{main}}_\ell(E)&=\sum_{q\mid N}v_\ell(\#\Phig_q), &
B^{\mathrm{tw}}_\ell(E)&=\sum_{q\mid N}w_\ell(E,q),
\end{align*}
and let $B^{\mathrm{III}}_\ell(E)$ be $B^{\mathrm{nc}}_\ell(E)$ with the terms $e_\ell(E,q)$ of the primes $q\ne\ell$ of type $\kod{III}$ or $\kod{III}^*$ left out.
So $B^{\mathrm{rule}}$ is the right side of (eq:rule), $B^{\mathrm{nc}}$ the same with no correction (``all types, no correction''), $B^{\mathrm{III}}$ the
``III-free'' form, $B^{\mathrm{main}}$ and $B^{\mathrm{tw}}$ the right sides of (eq:main) and (eq:twist). The right side of (eq:thm) is
$B^{\mathrm{thm}}_\ell(E)=\sum_{q\ne\ell}\bigl(e_\ell(E,q)+t_\ell(E,q)\bigr)$; it equals $B^{\mathrm{nc}}_\ell(E)$ when $\ell\nmid N$, and may be smaller when
$\ell\mid N$. For a bound $b$: a \emph{violation} is $v<b$; an \emph{equality} is $v=b$ with $b>0$; the \emph{excess} is $v-b$. Rates of
violation and of equality, and mean excesses, are always taken among the curves with $b>0$, as in the paper.
\end{convention}

\begin{convention}[populations and ranges]\label{conv:computation:populations}
\begin{enumerate}[label=(\alph*)]
\item A curve has \emph{CM} if its $j$-invariant is one of the thirteen rational CM $j$-invariants $0$, $1728$, $-3375$, $8000$, $-32768$, $54000$,
$287496$, $-884736$, $-12288000$, $16581375$, $-884736000$, $-147197952000$, $-262537412640768000$ (\Cref{ch:ellcurves}).
\item $E[\ell]$ is \emph{reducible} if and only if a label of \texttt{galrep} for $E$ begins with ``$\ell$B'' or ``$\ell$Cs'', the prime part
compared exactly (\Cref{conv:computation:labels}).
\item The \emph{base population at $\ell$} is the set of optimal curves without CM with $E[\ell]$ irreducible; the \emph{Eisenstein population at
$\ell$} is the set of optimal curves without CM with $E[\ell]$ reducible; the \emph{hypothesis population at $\ell$} is the base population
cut down by $\ell\nmid N$ and, at $\ell=3$, by ``image not $\Ns$'': the curves to which Theorem thm:main applies (Proposition prop:Sf).
\item The paper's ranges are the working set $N\le300000$, the hold-out $300000<N\le400000$, the extension $400000<N\le500000$ and the
out-of-table curves, $500000<N\le4\cdot10^6$. The book's ranges are $N<5000$ (the \emph{small working set}), $5000\le N<10000$ (the
\emph{small hold-out}) and their union; the split is for teaching only.
\end{enumerate}
Below conductor $10{,}000$ there are $38{,}042$ optimal curves (one per line of \texttt{degphi}), of which $189$ have CM. The base populations at
$\ell=3,5,7,11,13$ have $33{,}851$, $37{,}319$, $37{,}691$, $37{,}839$, $37{,}823$ curves, and the Eisenstein populations $4{,}002$, $534$, $162$, $14$, $30$; each
pair adds up to the $37{,}853$ curves without CM.
\end{convention}

Appendix A (\Cref{ch:prereq}) defines base rates, expectation, the inequalities of Markov and Chebyshev, hypothesis tests, hold-out sets and
pre-registration, and works the paper's calibration example in words; we use that vocabulary without repeating it. A reader with a
physics background knows the situation from experimental practice: the hold-out is a blinded signal region, the one-time test on it is the
unblinding, and the multiple-comparisons caveat of \Cref{rem:computation:multiple-comparisons} is the look-elsewhere effect.

Several earlier chapters send the reader here for the practical side of a computation: \Cref{ch:ellcurves} for the order of the output of
\texttt{ellinit} and the session at the end of its examples, \Cref{ch:reduction} for the formats of Cremona's files and the use of
\texttt{ellglobalred}, \Cref{ch:curves,ch:moddeg} for \texttt{ellmoddegree} and its memory needs, \Cref{ch:lfunctions} for installing and
calling PARI, \Cref{ch:groups} for the test \texttt{ns3}, and \Cref{ch:prereq} for the protocol and its counts. Each of these is met below.

\section{PARI/GP from zero}\label{sec:computation:pari}

This section is for the reader who has never used PARI/GP. It covers how to run it, the part of its language that the paper's code and
the book's scripts use, a first session on \curve{121a1} and \curve{726a1}, and the mistakes that produce wrong answers silently. In the paper
it is the background of the appendix.

\begin{remark}[running PARI]\label{rem:computation:install}
\begin{enumerate}[label=(\roman*)]
\item There are two ways to run PARI. The interpreter \texttt{gp} is distributed by the PARI group \cite{PARI} and packaged by the usual
Linux distributions; its manual \cite{PARIusers} documents every function used here. The Python package \texttt{cypari2} \cite{cypari2}
gives the same library inside Python: after \texttt{import cypari2; pari = cypari2.Pari()}, the call \texttt{pari('...')} evaluates any GP
expression and \texttt{pari('read("file.gp")')} loads a file of GP code. All outputs in this chapter come from the second route; the
results are those of the library and print exactly as \texttt{gp} prints them.
\item PARI has optional data packages. One of them, \texttt{elldata}, contains Cremona's tables inside PARI and is needed by
\texttt{ellsearch} and \texttt{ellidentify}. It is not installed in the book's environment, where \texttt{ellsearch(11)} fails with ``error
opening elldata file''. This is why the book reads the \texttt{ecdata} files directly (\Cref{sec:computation:ecdata}); a reader with
\texttt{elldata} can replace the file loops by \texttt{forell} or \texttt{ellsearch}.
\item PARI works on a stack. Two settings matter: \texttt{parisize}, the current size, and \texttt{parisizemax}, the size up to which PARI may
enlarge it when a computation needs more. In the book's installation a fresh \texttt{cypari2} session has \texttt{default(parisize)}
$=8{,}000{,}000$ and \texttt{default(parisizemax)} $=8{,}003{,}584$ bytes. In \texttt{gp} one writes \texttt{default(parisizemax, 10\^{}9)}; in
Python, \texttt{pari.allocatemem(s, smax)} sets both. A computation that needs more than \texttt{parisizemax} stops with ``the PARI stack
overflows'': an error, never a wrong answer. The function \texttt{ellmoddegree} is the only one here that needs more than the default
(\Cref{rem:computation:ellmoddegree}(iv)).
\end{enumerate}
\end{remark}

\begin{remark}[the part of the GP language used here]\label{rem:computation:gp-language}
Integers and rationals are exact, of arbitrary size. The letter \texttt{x} is PARI's default polynomial variable: never assign a value to
it, since \texttt{lfuneuler} returns rational functions in \texttt{x}. Vectors \texttt{[a, b, c]} are indexed from $1$; \texttt{\#v} is the
length and \texttt{v[2..\#v]} a slice; \texttt{v\textasciitilde} is the transpose (a column vector prints with a trailing
\texttt{\textasciitilde}). Matrices are written \texttt{[1, 2; 3, 4]}, with entry \texttt{M[i, j]} and column \texttt{M[, j]}; a $1\times2$ matrix
prints as \texttt{Mat([a, b])}. A function is defined by \texttt{f(a, \{b = 0\}) = ...}, the braces marking an optional argument with its
default; local variables are introduced by \texttt{my(...)}; a definition that spans several lines of a file is enclosed in braces
\texttt{\{ ... \}}; a comment runs from \texttt{\textbackslash\textbackslash} to the end of the line. Control flow: \texttt{if(c, a, b)},
\texttt{for(i = 1, n, ...)}, \texttt{foreach(v, s, ...)}, \texttt{while(c, ...)}, \texttt{next}, \texttt{return}. A \texttt{Map()} is a
dictionary (\texttt{mapput}, \texttt{mapisdefined}). Strings: \texttt{Str} concatenates, \texttt{strsplit(s, " ")} splits,
\texttt{readstr(file)} returns the vector of the lines of a file, and \texttt{eval} turns a string such as \texttt{"[0,-1,1,-10,-20]"}
into the vector it denotes (all available in 2.17.2). Arithmetic: \texttt{valuation(r, l)} is $v_\ell(r)$ for a rational $r$, and
\texttt{valuation(0, l)} is \texttt{+oo}, so that a test \texttt{vj < 0} is false when $j=0$; \texttt{factor} factors integers, rationals
and polynomials; \texttt{kronecker(D, n)} is the Kronecker symbol; \texttt{O(q\^{}k)} is the error term of $q$-adic precision $k$. Of the
functions that \Cref{ch:rings} names as counterparts of its algebra, \texttt{valuation} and \texttt{kronecker} are used here;
\texttt{znprimroot}, \texttt{ffgen} and \texttt{polisirreducible} are not needed by the paper.
\end{remark}

\begin{example}[a first session]\label{ex:computation:first-session}
On \curve{121a1} $=[1,1,1,-30,-76]$ and \curve{726a1} $=[1,1,0,-35,-51]$ (PARI/GP 2.17.2):
\begin{lstlisting}[language=GP]
? version()
[2, 17, 2]
? E = ellinit([1,1,1,-30,-76]);
? [E.a1, E.a2, E.a3, E.a4, E.a6]
[1, 1, 1, -30, -76]
? [E.b2, E.c4, E.c6, E.disc, E.j]
[5, 1441, 54703, -121, -24729001]
? ellan(E, 12)
[1, -1, 2, -1, 1, -2, 2, 3, 1, -1, 0, -2]
? ellap(E, 2)
-1
? E = ellinit([1,1,0,-35,-51]);
? elltors(E)
[2, [2], [[6, -3]]]
? ellisomat(E)[2]
[1, 2; 2, 1]
? ellisomat(ellinit([0,-1,1,-10,-20]))[2]
[1, 5, 5; 5, 1, 25; 5, 25, 1]
\end{lstlisting}
The member functions \texttt{E.a1}, \dots, \texttt{E.j} return the Weierstrass data of \Cref{ch:ellcurves}; for \curve{121a1}, $\Delta=-121=-11^2$ and
$j=-24729001$. The vector \texttt{ellan(E, 12)} lists $a_1,\dots,a_{12}$; $a_{11}=0$ because the reduction at $11$ is additive. For
\curve{726a1}, \texttt{elltors} returns the order of $E(\Q)_{\mathrm{tors}}$, its structure and a generator: $\Z/2$, generated by $(6,-3)$.
The second component of \texttt{ellisomat} is the matrix of isogeny degrees between the curves of the class: the class of \curve{726a1}
has two curves joined by a $2$-isogeny, and that of \curve{11a1} three curves joined by $5$-isogenies, with the degree-$25$ isogeny between
the second and the third. These are the matrices of the \texttt{allisog} file (\Cref{def:computation:ecdata-files}).
\end{example}

\begin{warning}[pitfalls]\label{warn:computation:pitfalls}
\begin{enumerate}[label=(\alph*)]
\item \texttt{ellinit} accepts any Weierstrass model. Local data (\texttt{elllocalred}, \texttt{ellglobalred}) are always those of a minimal
model, and \texttt{ellglobalred} reports the change of variables to one; \texttt{ellminimalmodel} returns the minimal model itself. The
member functions, on the other hand, are those of the model given.
\item Never assign to \texttt{x} (\Cref{rem:computation:gp-language}).
\item A $q$-adic object carries its precision \texttt{O(q\^{}k)}; results are correct only to that precision.
\item A stack overflow is an error, not a wrong answer: raise \texttt{parisizemax} and rerun.
\item GP indexes from $1$ and Python from $0$. In \texttt{cypari2}, the GP entry \texttt{R[5][i][2]} is \texttt{R[4][i-1][1]}, and the GP
matrix entry \texttt{R[4][i, 1]} is \texttt{R[3][i-1, 0]}. \texttt{cypari2} does not accept Python slices of a PARI matrix (the expression
\texttt{R[3][:,0]} raises a \texttt{TypeError}), so whole columns are best taken in GP, for instance with \texttt{pari('(g) -> g[4][,1]\textasciitilde')(R)}.
\item \texttt{factor} of a polynomial with integer coefficients returns its irreducible factors in $\Z[x]$, primitive and with positive
leading coefficient, and drops the content, sign included: \texttt{factor(3*x\^{}4 + 6)} is \texttt{Mat([x\^{}4 + 2, 1])}
(\Cref{ex:computation:mod3-image}). This is why the paper's \texttt{ns3} does not see the leading coefficient $3$ of $\psi_3$.
\item A file of GP code must be loaded with \texttt{read}. Pasted into \texttt{pari('...')} as one string, the paper's file
\texttt{euler\_rule.gp} does not produce two functions: in our test the definition of \texttt{ns3} became part of the body of
\texttt{rule}, and \texttt{rule(E, 3)} returned a function instead of a pair. Reading the file is the reliable route.
\item \texttt{mfsturm} depends on its argument. On a pair $[N,k]$ it returns Sturm's bound $k\mu(N)/12$, $\mu(N)=[\SL_2(\Z):\Gamma_0(N)]$
(\Cref{ch:modforms}): \texttt{mfsturm([539, 2])} is $112$. On a space created by \texttt{mfinit} it returns a smaller number adapted to that
space: \texttt{mfsturm(mfinit([539, 2]))} is $88$ and, for the cuspidal space, \texttt{mfsturm(mfinit([539, 2], 1))} is $50$. The smaller numbers
are not valid bounds for the integrality of an arbitrary rational combination of $q$-expansions; used as truncation points in the
congruence-module computation of \Cref{ch:moddeg} at level $539$, the bound $50$ gives the wrong modules ($1440$ and $2520$ instead of
$288$ and $504$, as found while writing \Cref{ch:congruence}). For congruences and integrality, use $k\mu(N)/12$.
\end{enumerate}
\end{warning}
\section{Invariants of a curve in PARI}\label{sec:computation:invariants}

This section goes through the PARI functions that compute the local invariants the paper uses: the local data of \texttt{ellglobalred}, the
Euler factors of the symmetric square, the Tate parameter, the $3$-division polynomial, quadratic twists and the modular degree. For each it
says what the output is, prints an exact session, and checks a lemma of the paper against PARI on every optimal curve below conductor
$10{,}000$. In the paper this is the content of the sentence of Section sec:data that lists what was recomputed with PARI, and of the checks
claimed in Lemma lem:euler, Lemma lem:twist and the appendix.

\begin{definition}[the output of \texttt{ellglobalred} and \texttt{elllocalred}]\label{def:computation:ellglobalred}
For an elliptic curve $E$ over $\Q$ given by any model, \texttt{elllocalred(E, p)} is the vector $[f_p,k_p,[u,r,s,t],c_p]$: the conductor
exponent at $p$, PARI's Kodaira code $k_p$ (\Cref{conv:computation:code-idioms}), the change of variables to a model minimal at $p$, and the
Tamagawa number $c_p=\#\Phi_p(\F_p)$. The vector \texttt{R = ellglobalred(E)} is $[N,[u,r,s,t],\prod_pc_p,F,L]$: the conductor; the change of variables to a
global minimal model; the product of the Tamagawa numbers; the factorisation matrix $F$ of $N$, whose first column \texttt{R[4][,1]} lists the bad primes in
increasing order; and the vector $L$ with $L[i]$ equal to \texttt{elllocalred(E, p\_i)} for the $i$-th bad prime, except that over $\Q$ the
local change of variables is replaced by $0$ \cite{PARIusers}. So the paper's \texttt{P = R[4][,1]} is the column of bad primes and its
\texttt{K[i][2]}, with \texttt{K = R[5]}, the Kodaira code at the $i$-th of them.
\end{definition}

For \curve{726a1} (PARI/GP 2.17.2):
\begin{lstlisting}[language=GP]
? E = ellinit([1,1,0,-35,-51]);
? R = ellglobalred(E)
[726, [1, 0, 0, 0], 4, [2, 1; 3, 1; 11, 2], [[1, 10, 0, 2], [1, 7, 0, 1], [2, 3, 0, 2]]]
? R[4][,1]
[2, 3, 11]~
? R[5][3][2]
3
? elllocalred(E, 11)
[2, 3, [1, 0, 0, 0], 2]
? elllocalred(ellinit([0,1,0,4,4]), 2)
[2, -4, [1, 0, 0, 0], 3]
\end{lstlisting}
The model is already minimal (the change of variables is the identity $[1,0,0,0]$), $N=2\cdot3\cdot11^2$, the types are $\kod{I}_6$, $\kod{I}_3$, $\kod{III}$ (codes
$10$, $7$, $3$) with $c_2=2$, $c_3=1$, $c_{11}=2$ and product $4$. For \curve{20a1} $=[0,1,0,4,4]$ at $2$ the code $-4$ is $\kod{IV}^*$, with $f_2=2$ and
$c_2=3$. The mathematics (Tate's algorithm, why $c_q$ divides $\#\Phig_q$ and can be smaller) is in \Cref{ch:reduction}; this is the data structure only.
Cremona's file \texttt{curvedata} stores the same bad primes, codes, Tamagawa numbers and conductor exponents, and they agree with
\texttt{ellglobalred} on every curve below conductor $10{,}000$ (\Cref{ex:computation:curvedata}).

\begin{convention}[Kodaira codes in code]\label{conv:computation:code-idioms}
PARI's code $k$ is $1$ for $\kod{I}_0$; $n+4$ for $\kod{I}_n$ and $-(n+4)$ for $\kod{I}_n^*$, \emph{for $n\ge1$ only}; $-1$ for $\kod{I}_0^*$; $2,3,4$ for
$\kod{II},\kod{III},\kod{IV}$ and $-2,-3,-4$ for $\kod{II}^*,\kod{III}^*,\kod{IV}^*$ (\Cref{ch:reduction}). The paper's appendix states the codes ``$n+4$ for
$\mathrm{I}_n$ and $-(n+4)$ for $\mathrm{I}_n^*$ when $n\ge1$''; read without the last clause, $n+4$ at $n=0$ would give $4$, the code of $\kod{IV}$. The
tests used in the paper's code and in the book's scripts:
\begin{center}\small
\begin{tabular}{lll}
\toprule
test on $k$ & types selected & use\\
\midrule
\texttt{k >= 5} & $\kod{I}_n$, $n=k-4\ge1$ & multiplicative: divide by the naive factor; $\#\Phig_q=n$\\
\texttt{k <= -5} & $\kod{I}_n^*$, $n=-k-4\ge1$ & $t_\ell=v_\ell(n)$ when $q$ is odd\\
\texttt{k == -1} & $\kod{I}_0^*$ & $\#\Phig_q=4$\\
\texttt{abs(k) == 3} & $\kod{III}$, $\kod{III}^*$ & left out unless \texttt{iii} is set\\
\texttt{k == 4 || k == 2} & $\kod{IV}$, $\kod{II}$ & $w_3=1$ at $q=\ell=3$\\
\texttt{k == -4 || k == -2} & $\kod{IV}^*$, $\kod{II}^*$ & $w_3=2$ at $q=\ell=3$\\
\texttt{k == -3} & $\kod{III}^*$ & $w_3=1$ at $q=\ell=3$\\
\bottomrule
\end{tabular}
\end{center}
\end{convention}

\begin{remark}[the symmetric square in PARI]\label{rem:computation:lfunsympow}
\texttt{S = lfunsympow(E, 2)} is PARI's $L$-function object for $L(\Sym f,s)$, in the normalisation of \Cref{ch:lfunctions}, with the true local factors
at every prime. \texttt{lfunparams(S)} returns the conductor of $L(\Sym f,s)$, the number $k$ of the functional equation $s\leftrightarrow k-s$, and
the gamma shifts. \texttt{lfuneuler(S, q)} is the Euler factor at $q$ as a rational function of $X=q^{-s}$ (printed in the variable
\texttt{x}): it is $1/P_q(X)$. Hence $P_q(q^{-2})=L_q(A_f,1)^{-1}$ is the exact rational \texttt{subst(1/lfuneuler(S, q), x, 1/q\^{}2)}, and
$e_\ell(E,q)$ is the $\ell$-adic valuation of its quotient by $1-a_q^2q^{-2}$ (Definition def:terms). For \curve{121a1} (PARI/GP 2.17.2):
\begin{lstlisting}[language=GP]
? S = lfunsympow(ellinit([1,1,1,-30,-76]), 2);
? lfunparams(S)
[121, 3, [0, 0, 1]]
? lfuneuler(S, 11)
1/(11*x + 1)
? lfuneuler(S, 2)
1/(-8*x^3 - 2*x^2 + x + 1)
? subst(1/lfuneuler(S, 11), x, 1/11^2)
12/11
\end{lstlisting}
The conductor of $L(\Sym f,s)$ is $121$, the functional equation relates $s$ and $3-s$, and the gamma factor has shifts $0,0,1$. At the additive
prime $11$ (type $\kod{II}$, $11\equiv2\pmod3$) the factor is $P_{11}(X)=1+11X$, the row $\kod{II}$ of the paper's table with $\chi_{-3}(11)=-1$, and
$P_{11}(11^{-2})=12/11$, so $e_3(E,11)=1$. At the good prime $2$, with $a_2=-1$, $P_2(X)=1+X-2X^2-8X^3=(1-2X)(1+3X+4X^2)$, the factor
$(1-\alpha^2X)(1-2X)(1-\beta^2X)$ with $\alpha+\beta=-1$, $\alpha\beta=2$. The bad factors are computed when \texttt{S} is created; creating
\texttt{S} and reading all bad factors took about $0.08$ milliseconds per curve on the first $3{,}000$ optimal curves of the tables. The
general constructor \texttt{lfuncreate} is not needed. \texttt{lfun(S, 2)} is the value $L(\Sym f,2)$ that \texttt{ellmoddegree} uses.
\end{remark}

\begin{example}[the omitted factors by hand, and the check of Lemma lem:euler]\label{ex:computation:euler-by-hand}
For a bad prime $q\ge5$, Lemma lem:euler reads $L_q(A_f,1)^{-1}$ off the Kodaira type (\Cref{ch:lfunctions}). The function \texttt{euler\_by\_type} of the
script \texttt{sessions.gp} (\Cref{con:computation:sessions}) implements the table: $1-q^{-2}$ for $\kod{I}_n$ and $\kod{I}_n^*$; $1-\chi_{-3}(q)/q$ for $\kod{II}$, $\kod{II}^*$,
$\kod{IV}$, $\kod{IV}^*$; $1-\chi_{-4}(q)/q$ for $\kod{III}$, $\kod{III}^*$; and $(q-1)\bigl((q+1)^2-a'^2\bigr)/q^3$ for $\kod{I}_0^*$, where $a'$ is $a_q$ of the twist by
$q^*$, computed as \texttt{ellap(ellinit(elltwist(E, qs)), q)} with \texttt{qs} $=q^*$. The function \texttt{eulercheck(E)} returns the bad primes
$q\ge5$ at which this differs from PARI's factor. For \curve{539d1} $=[0,1,1,-16,-66]$, of type $\kod{I}_0^*$ at $7$, with $a'=-2$ (PARI/GP 2.17.2):
\begin{lstlisting}[language=GP]
? E = ellinit([0,1,1,-16,-66]);
? euler_by_type(E, 7)
360/343
? subst(1/lfuneuler(lfunsympow(E, 2), 7), x, 1/7^2)
360/343
? eulercheck(E)
[]
\end{lstlisting}
So $e_3(E,7)=v_3(360/343)=2$: $(7-1)\bigl(64-4\bigr)/7^3=360/343$. \emph{Check.} On all $67{,}745$ pairs (optimal curve of conductor below $10{,}000$, bad prime
$q\ge5$) the two values agree: $58{,}860$ pairs of type $\kod{I}_n$, $2{,}709$ of type $\kod{I}_n^*$, $3{,}588$ of types $\kod{II}$, $\kod{II}^*$, $\kod{IV}$, $\kod{IV}^*$,
$1{,}676$ of types $\kod{III}$, $\kod{III}^*$ and $912$ of type $\kod{I}_0^*$, with no mismatch. This is the reader's version of the paper's sentence ``In each
case the values were confirmed against PARI's \texttt{lfunsympow} on every curve of the tables'' (proof of Lemma lem:euler); \Cref{ch:lfunctions}
reports the same check over all curves.
\end{example}

\begin{example}[the Tate parameter]\label{ex:computation:tate-parameter}
\texttt{Eq = ellinit(E, O(q\^{}k))} is $E$ over $\Q_q$ with $q$-adic precision $k$, and when $v_q(j)<0$ its member \texttt{Eq.tate} has as third
component the Tate parameter $q_E\in\Q_q$, with $v_q(q_E)=-v_q(j)$ (\Cref{ch:reduction}). The function \texttt{tatedata(E, q, l)} of
\texttt{sessions.gp} returns $[n,q_E,\bar u,k]$: $n=v_q(q_E)$; $q_E$ to precision $q^{30}$; $\bar u$, the reduction modulo $q$ of the unit part
$u=q_Eq^{-n}$, as an integer in $[0,q-1]$; and the largest $k\le v_\ell(q-1)$ such that $\bar u$ is an $\ell^k$-th power in $\F_q^\times$, found by the
test $\bar u^{(q-1)/\ell^k}=1$. In the notation of the paper and of \Cref{ch:blochkato}, $k=v_\ell(c(u))$, the valuation of the class of $\bar u$; we write
$c(u)$ and never $c_q$, which is a Tamagawa number. Sessions (PARI/GP 2.17.2):
\begin{lstlisting}[language=GP]
? E = ellinit([0,-1,1,-10,-20]);
? ellinit(E, O(11^8)).tate[3]
10*11^5 + 2*11^6 + 6*11^7 + 6*11^8 + 5*11^9 + 4*11^10 + 4*11^11 + 11^12 + O(11^13)
? T = tatedata(E, 11, 5); [T[1], T[3], T[4]]
[5, 10, 1]
? T = tatedata(ellinit([1,1,0,4,11]), 7, 3); [T[1], T[3], T[4]]
[3, 5, 0]
? T = tatedata(ellinit([1,0,1,170,-3237]), 7, 3); [T[1], T[3], T[4]]
[3, 5, 0]
? tatedata(ellinit([1,1,0,4,11]), 7, 3)[2] == tatedata(ellinit([1,0,1,170,-3237]), 7, 3)[2]
1
? T = tatedata(ellinit([1,0,1,-31,117]), 7, 3); [T[1], T[3], T[4]]
[3, 6, 1]
? T = tatedata(ellinit([0,-1,1,30,-335]), 19, 3); [T[1], T[3], T[4]]
[3, 8, 1]
? ellinit(ellinit([1,1,0,4,11]), 7)
[]
\end{lstlisting}
For \curve{11a1} at $11$: $n=5$, $q_E=10\cdot11^5+2\cdot11^6+\cdots$, $\bar u=10\equiv-1$, a fifth power in $\F_{11}^\times$ (the fifth powers are $\pm1$), so
$k=1$ at $\ell=5$. For \curve{77c1} $=[1,1,0,4,11]$ at $7$ (type $\kod{I}_3$, $a_7=-1$, nonsplit): $n=3$, $\bar u=5$, not a cube in $\F_7^\times$ (the cubes
are $\pm1$), $k=0$. Its twist \curve{539c1} by $\chi_{-7}$ (type $\kod{I}_3^*$ at $7$) has the same $j$, hence the same Tate parameter: PARI returns
the same series. For \curve{791a1} $=[1,0,1,-31,117]$ at $7$, $\bar u=6\equiv-1$ is a cube; for \curve{1577a1} $=[0,-1,1,30,-335]$ at $19$, $\bar u=8=2^3$ is a cube
but not a ninth power, so $k=1<2=v_3(18)$. These are the examples of \Cref{ch:galcoh,ch:blochkato}. Three warnings. First, \texttt{elltatepairing} is the
Tate pairing on torsion points and has nothing to do with the Tate parameter. Second, \texttt{ellinit(E, q)} with an integer prime $q$
reduces $E$ modulo $q$: at the bad prime $7$ of \curve{77c1} it returns \texttt{[]}, and at a good prime the curve over $\F_q$, never the curve over
$\Q_q$, which needs the precision \texttt{O(q\^{}k)}. Third, at $q=2$ PARI 2.17.2 returns a wrong parameter when $v_2(j)=-1$
(\Cref{ch:reduction}: the output is $q_E^5$ instead of $q_E$; for \curve{208a1} $=[0,-1,0,8,-16]$, of type $\kod{I}_1^*$ at $2$, PARI returns a parameter of
valuation $5$, equal to the fifth power of the one obtained from $j$); there one inverts the series $j=\mathsf q^{-1}+744+196884\mathsf q+\cdots$ of \Cref{ch:reduction} instead. The paper's
Tamagawa test never needs $q=2$, since $q\equiv1\pmod\ell$ fails for $q=2$ and every odd $\ell$; the $t_\ell$ terms use only $v_q(j)$.
\end{example}

\begin{example}[the mod $3$ image from the $3$-division polynomial]\label{ex:computation:mod3-image}
\texttt{elldivpol(E, 3)} is the $3$-division polynomial $\psi_3$, of degree $4$ and leading coefficient $3$, whose roots are the $x$-coordinates of the
eight points of order $3$. For \curve{726a1} (PARI/GP 2.17.2):
\begin{lstlisting}[language=GP]
? E = ellinit([1,1,0,-35,-51]);
? elldivpol(E, 3)
3*x^4 + 5*x^3 - 210*x^2 - 612*x - 1480
? factor(elldivpol(E, 3))
[x^2 - x - 74, 1; 3*x^2 + 8*x + 20, 1]
? ns3(E)
1
? factor(3*x^4 + 6)
Mat([x^4 + 2, 1])
? ns3(ellinit([0,-1,1,-7,10]))
1
\end{lstlisting}
The pattern $2+2$ is that of image $N_s(\F_3)$ when $E[3]$ is irreducible (Proposition prop:ns; \Cref{ch:galrep} has the table of patterns and
images), and the paper's \texttt{ns3(E)} returns $1$. \texttt{polgalois} requires an irreducible polynomial: on this $\psi_3$ it stops with ``not
an irreducible polynomial'', so it can be used only when $\psi_3$ is irreducible. \emph{Check.} Among the $37{,}853$ optimal curves without CM below
conductor $10{,}000$, \texttt{ns3(E)} $=1$ exactly for the $90$ curves whose \texttt{galrep} entry contains \texttt{3Ns}, with no mismatch either way.
Among the $189$ optimal CM curves, \texttt{ns3} is $1$ for $13$: \curve{121b1}, \curve{256a1}, \curve{256d1}, \curve{1089e1}, \curve{1936f1}, \curve{2304c1},
\curve{2304k1}, \curve{3025a1}, \curve{5929a1}, \curve{6400f1}, \curve{6400r1}, \curve{7744a1}, \curve{7744s1}, all with $j=-32768$ or $8000$ (the last session line is
\curve{121b1}). Their image is $N_s(\F_3)$, but \texttt{galrep} gives them no label: Cremona's file omits the images forced by complex
multiplication (\Cref{ch:galrep}). This is why CM is decided by $j$ (\Cref{conv:computation:populations}).
\end{example}

\begin{example}[quadratic twists, and the check of Lemma lem:twist]\label{ex:computation:twists}
\texttt{elltwist(E, d)}, for a discriminant $d$, is a model of the quadratic twist of $E$ by $\Q(\sqrt d)$, not necessarily minimal;
\texttt{ellminimalmodel} returns Cremona's reduced minimal model, which can be looked up in \texttt{allcurves}. The function
\texttt{twistcheck(E)} of \texttt{sessions.gp} lists, for each odd bad prime $q$, the code of $E$ at $q$, the code at $q$ of the twist $E^{(q)}$ by
$q^*$, and the conductor of $E^{(q)}$ (PARI/GP 2.17.2):
\begin{lstlisting}[language=GP]
? twistcheck(ellinit([1,1,1,-30,-76]))
[[11, 2, -4, 121]]
? twistcheck(ellinit([1,0,1,170,-3237]))
[[7, -7, 7, 77], [11, 6, -6, 5929]]
? twistcheck(ellinit([0,0,0,1568,-72716]))
[[5, 5, -5, 4900], [7, -9, 9, 140]]
? E = ellminimalmodel(ellinit(elltwist(ellinit([1,1,1,-30,-76]), -11))); [E.a1, E.a2, E.a3, E.a4, E.a6]
[1, 1, 0, -3632, 82757]
\end{lstlisting}
\curve{121a1} ($\kod{II}$ at $11$) twists to a curve of type $\kod{IV}^*$ and the same conductor, \curve{121c2} $=[1,1,0,-3632,82757]$; \curve{539c1} goes to
\curve{77c} at $7$ ($\kod{I}_3^*\to\kod{I}_3$) and to conductor $5929=7^2\cdot11^2$ at $11$ ($\kod{I}_2\to\kod{I}_2^*$); \curve{980i1} goes to conductor $4900$ at $5$ and to
\curve{140b} at $7$ ($\kod{I}_5^*\to\kod{I}_5$). \emph{Check.} For every optimal curve below conductor $10{,}000$ and every odd bad prime $q$ we compared the code and the
conductor exponent of $E^{(q)}$ at $q$ with Lemma lem:twist(ii): $\kod{I}_n\leftrightarrow\kod{I}_n^*$ and $\kod{I}_0^*\to\kod{I}_0$ at every odd $q$, and the exchanges
$\kod{II}\leftrightarrow\kod{IV}^*$, $\kod{III}\leftrightarrow\kod{III}^*$, $\kod{IV}\leftrightarrow\kod{II}^*$ at $q\ge5$. All $85{,}848$ pairs agree ($72{,}003$ of type
$\kod{I}_n$, $6{,}348$ of type $\kod{I}_n^*$, $2{,}233$ of type $\kod{I}_0^*$, $5{,}264$ potentially good at $q\ge5$), and in every pair the code and the conductor
exponent at every other bad prime, $2$ included, are unchanged, which is Lemma lem:twist(i). At $q=3$ the same three exchanges, with $f_3$
unchanged, hold in all $5{,}122$ pairs where $E$ has potentially good additive reduction at $3$ of one of these six types; \Cref{ch:reduction}
proves this. The paper's own script \texttt{twist\_lemma\_check.gp} checked $4{,}577$ random pairs (the appendix).
\end{example}

\begin{remark}[the modular degree in PARI]\label{rem:computation:ellmoddegree}
\begin{enumerate}[label=(\roman*)]
\item \texttt{ellmoddegree(E)} returns the modular degree of the given curve divided by the square of its Manin constant \cite{PARIusers}: the
degree $\deg\varphi$ of the minimal parametrisation $X_0(N)\to E$ divided by $c_\varphi^2$ (\Cref{ch:moddeg}). On the optimal curve, where
$c_E=1$ in all of Cremona's tables, it is $\degphi$; on another curve of the class it can be a fraction. For the class \curve{11a}:
\begin{lstlisting}[language=GP]
? [ellmoddegree(ellinit([0,-1,1,-10,-20])), ellmoddegree(ellinit([0,-1,1,-7820,-263580])), ellmoddegree(ellinit([0,-1,1,0,0]))]
[1, 5, 1/5]
? W = ellweilcurve(ellinit([0,-1,1,0,0])); W[2]
[[1, 1/5], [1, 1], [5, 1]]
? apply(e -> [e.a4, e.a6], W[1])
[[0, 0], [-10, -20], [-7820, -263580]]
? W = ellweilcurve(ellinit([1,-1,1,-96608,-11533373])); W[2]
[[3, 1], [1, 1], [6, 1], [2, 1]]
? apply(e -> [e.a4, e.a6], W[1])
[[-96608, -11533373], [-1568, -4669], [-95528, -11804669], [6112, -41533]]
? [ellmoddegree(ellinit([1,-1,1,-96608,-11533373])), ellmoddegree(ellinit([1,-1,1,-1568,-4669]))]
[5184, 1728]
\end{lstlisting}
\item \texttt{ellweilcurve(E)} returns $[vE,vS]$: minimal models of the curves of the class (in the order of \texttt{ellisomat}) and, for each,
the Smith invariants of its lattice of modular symbols inside the lattice of the newform; the strong Weil (optimal) curve is the one with
invariants $[c_E,c_E]$, $c_E$ its Manin constant \cite{PARIusers}. Called on \curve{11a3} it lists \curve{11a3}, \curve{11a1}, \curve{11a2}, and the entry
$[1,1]$ marks \curve{11a1}; called on \curve{990h1} it lists \curve{990h1}, \curve{990h3}, \curve{990h2}, \curve{990h4}, and $[1,1]$ marks \curve{990h3}. This is how PARI
identifies the optimal curve for moderate conductors.
\item The method is Watkins's \cite{Watkins02}: compute $L(\Sym f,2)$ with the true bad factors, deduce the Petersson norm, and divide by the
area of the period lattice (\Cref{ch:moddeg,ch:lfunctions}). The \texttt{lfun} machinery prints ``\texttt{\_lfuninit\_worker: Warning:
increasing stack size}'' when it enlarges the stack.
\item \emph{Cost and memory} (PARI/GP 2.17.2). A random sample of $50$ optimal curves below conductor $10{,}000$ took $0.135$ seconds per curve on a
two-core machine that was running the full check of \Cref{ex:computation:moddeg-check} at the same time. With the default stack
(\Cref{rem:computation:install}) \texttt{ellmoddegree} succeeds on \curve{9999m1} ($\degphi=5632$) but overflows on the curves of largest degree below
$10{,}000$: \curve{9510e1} ($\degphi=100{,}254{,}000$) succeeds with a $16$ MB stack, \curve{9690m1} ($68{,}217{,}600$) needs more than $16$ MB and succeeds
with $32$ MB, and \curve{7910b1} ($31{,}272{,}384$) overflows the default and succeeds with $32$ MB. This agrees with the paper's appendix: ``The
function \texttt{ellmoddegree} needs about $32$ MB of stack with PARI 2.17 for conductors in the thousands (with PARI 2.15 and larger
conductors, \texttt{parisizemax} of a few hundred megabytes).''
\item \emph{Checks.} \texttt{ellmoddegree} agrees with \texttt{degphi} on all $38{,}042$ optimal curves below $10{,}000$
(\Cref{ex:computation:moddeg-check}); and on all $2{,}650$ non-optimal curves of conductor below $1000$, \texttt{ellmoddegree(E)} equals the degree
of \texttt{alldegphi} divided by the square of the Manin constant of \texttt{opt\_man}, the value being a fraction for $33$ of them (\curve{11a3},
\curve{14a4}, \curve{14a6}, \curve{15a3}, \dots).
\end{enumerate}
\end{remark}

\begin{paperlink}\label{pl:computation:invariants}
Section sec:data: ``The Kodaira types and discriminant exponents at the bad primes were recomputed with PARI/GP 2.15 \cite{PARI} from the
$a$-invariants, the Euler factors of $\Sym f$ at every bad prime were computed with \texttt{lfunsympow} and \texttt{lfuneuler}, the Tate
parameters with \texttt{ellinit} over $\Q_q$, and fifteen modular degrees spanning the configurations were recomputed with
\texttt{ellmoddegree}, all agreeing with the tables.'' Each of these computations is one of the items above
(\Cref{def:computation:ellglobalred,rem:computation:lfunsympow,ex:computation:tate-parameter,rem:computation:ellmoddegree}); below
conductor $10{,}000$ every one of them has been run on every optimal curve.
\end{paperlink}
\section{Cremona's tables and the \texttt{ecdata} files}\label{sec:computation:ecdata}

The paper's data are Cremona's tables of elliptic curves over $\Q$ \cite{Cremona97}, in the form of the \texttt{ecdata} repository
\cite{ecdata}. This section describes the files used in this book, shows how to read them in GP and in Python, records what the files say
about optimality, and fixes the labels. In the paper it is the background of the first paragraph of Section sec:data and of Section 1.4.

\begin{definition}[the \texttt{ecdata} files]\label{def:computation:ecdata-files}
The repository consists of plain text files, one per range of $10{,}000$ conductors, grouped in directories by kind; each line describes one
curve or one isogeny class, with fields separated by single spaces, and $a$-invariants written \texttt{[a1,a2,a3,a4,a6]} without spaces. A
Cremona label $N$, class, number is written in three fields (\texttt{11 a 1}) or as one word (\texttt{11a1}). The six files used for the counts of
this chapter, for conductors $0$ to $9{,}999$, have the following formats (read off the files and the repository's description
\texttt{docs/file-format.txt}):
\begin{center}\small
\begin{tabular}{lrlp{8.2cm}}
\toprule
file & lines & one line per & fields\\
\midrule
\texttt{allcurves} & $64{,}687$ & curve & $N$, class, number, $[a_i]$, rank, $\#E(\Q)_{\mathrm{tors}}$\\
\texttt{allbsd} & $64{,}687$ & curve & as \texttt{allcurves}, then $\prod_pc_p$, $\Omega$, $L^{(r)}(E,1)/r!$, regulator, analytic order of Sha\\
\texttt{allisog} & $38{,}042$ & class & $N$, class, $1$, $[a_i]$ of the optimal curve, the list of $[a_i]$ of the class, the isogeny matrix\\
\texttt{alllabels} & $64{,}687$ & curve & Cremona's $N$, class, number, then the LMFDB's $N$, class, number\\
\texttt{degphi} & $38{,}042$ & class & $N$, class, $1$, $\degphi$, the set of primes dividing it, $[a_i]$ of the optimal curve\\
\texttt{galrep} & $64{,}687$ & curve & label, then the labels of the non-surjective mod $\ell$ images\\
\bottomrule
\end{tabular}
\end{center}
Sample lines, in order from \texttt{allcurves}, \texttt{allbsd}, \texttt{allisog}, \texttt{alllabels}, \texttt{degphi} (three lines) and
\texttt{galrep} (two lines):
\begin{lstlisting}
11 a 1 [0,-1,1,-10,-20] 0 5
11 a 1 [0,-1,1,-10,-20] 0 5 5 1.26920930427955 0.253841860855911 1.00000000000000 1
11 a 1 [0,-1,1,-10,-20] [[0,-1,1,-10,-20],[0,-1,1,-7820,-263580],[0,-1,1,0,0]] [[1,5,5],[5,1,25],[5,25,1]]
11 a 2 11 a 1
11 a 1 1 {} [0,-1,1,-10,-20]
726 a 1 144 {2,3} [1,1,0,-35,-51]
990 h 1 1728 {2,3} [1,-1,1,-1568,-4669]
11a1 5Cs.1.1
726a1 2B 3Ns
\end{lstlisting}
So \curve{11a1} has rank $0$, torsion of order $5$, $\prod_pc_p=5$, $\Omega=1.2692\dots$ and $L(E,1)=\Omega/5$; its class has three curves; its
modular degree is $1$; and its mod $5$ image is \texttt{5Cs.1.1}. A \texttt{galrep} line lists an image only for the primes $\ell$ at which it is
not $\GL_2(\F_\ell)$: $20{,}930$ of the $64{,}687$ lines have no label at all. In the \texttt{degphi} and \texttt{allisog} files the third field is always
$1$ and the equation is that of the optimal curve, which is curve number $1$ of its class except in one class (\Cref{rem:computation:optimality-data}).
\end{definition}

The repository holds more: among the directories at commit \texttt{25cec5e} are \texttt{allgens} (generators), \texttt{aplist}
($a_p$ for $p<100$), \texttt{intpts} (integral points), \texttt{alldegphi} (one line per curve: the degree of the minimal parametrisation of
\emph{every} curve of the class), \texttt{opt\_man} (one line per curve: a flag marking the optimal curve, and the Manin constant),
\texttt{curvedata} (one line per curve, with many fields, among them the bad primes, the Tamagawa numbers, the Kodaira symbols and the
conductor exponents), \texttt{classdata} (one line per class: the isogeny matrix and the first coefficients of the newform), and \texttt{2adic}. The six files of the table do not contain the Kodaira symbols
or the local Tamagawa numbers (only $\prod_pc_p$, in \texttt{allbsd}); \texttt{curvedata} does. So the paper's description of the
repository (``which contain every elliptic curve over $\Q$ of conductor below $500000$ with its Kodaira symbols, Tamagawa numbers, modular
degree, isogeny class, torsion and the images of the mod $\ell$ Galois representations'') is accurate for the repository as a whole.

\begin{example}[cross-checks against the further files]\label{ex:computation:curvedata}
At commit \texttt{25cec5e}, for conductors $0$ to $9{,}999$:
\begin{enumerate}[label=(\alph*)]
\item the bad primes, the Kodaira codes (stored in PARI's coding), the Tamagawa numbers and the conductor exponents of \texttt{curvedata}
agree with \texttt{ellglobalred} on all $64{,}687$ curves;
\item \texttt{opt\_man} flags exactly one curve in each of the $38{,}042$ classes as optimal, always the curve whose equation is in
\texttt{degphi}, and gives it Manin constant $1$; over all curves the Manin constants are $1$ ($64{,}592$ curves), $2$ ($70$), $3$ ($22$), $4$ ($2$)
and $5$ ($1$, the curve \curve{11a3});
\item the degree that \texttt{alldegphi} gives for the optimal curve of each class equals the degree of \texttt{degphi}, in all $38{,}042$ classes.
\end{enumerate}
\end{example}

\begin{example}[reading the files]\label{ex:computation:reading-ecdata}
In GP, \texttt{readstr(file)} returns the vector of lines, \texttt{strsplit(s, " ")} the fields of a line, and \texttt{eval} the vector
written in a field (PARI/GP 2.17.2):
\begin{lstlisting}[language=GP]
? Vec(strsplit("11 a 1 [0,-1,1,-10,-20] 0 5", " "))
["11", "a", "1", "[0,-1,1,-10,-20]", "0", "5"]
? eval("[0,-1,1,-10,-20]")
[0, -1, 1, -10, -20]
\end{lstlisting}
The function \texttt{readgal} of \Cref{con:computation:table-small} builds a \texttt{Map} from labels to lists of image labels, and its loop
\texttt{tablerow} runs over \texttt{degphi}. In Python the same is \texttt{line.split()} and \texttt{ast.literal\_eval} (or \texttt{eval}, on
trusted data). One caution: the label of the optimal curve must not be assembled from the three label fields of a \texttt{degphi} line, since
the number field is $1$ even for the class \curve{990h}, whose optimal curve is \curve{990h3}. The book's Python script looks the equation up in
\texttt{allcurves}; the GP script treats the one exception below $10{,}000$ explicitly.
\end{example}

\begin{remark}[optimality in the files]\label{rem:computation:optimality-data}
\begin{enumerate}[label=(\roman*)]
\item In Cremona's labelling, curve number $1$ of a class is the optimal curve with exactly one exception below conductor $10{,}000$: of the
$38{,}042$ equations of \texttt{degphi}, exactly one is not curve number $1$ of \texttt{allcurves}, namely \curve{990h3} $=[1,-1,1,-1568,-4669]$
(\Cref{ch:ellcurves}). The consequence for computation: \texttt{ellmoddegree} gives $1728$ on \curve{990h3} and $5184$ on \curve{990h1}, and
\texttt{rule} gives \texttt{[3, 3]} and \texttt{[4, 3]} at $\ell=3$ (\Cref{warn:computation:rule-pitfalls}).
\item The paper: ``Optimality is determined in the tables for every class of conductor at most $400000$ and for $359{,}009$ further classes
up to $500000$'' (Section sec:data). What ``determined'' means, and how Cremona certifies it, is in \Cref{ch:moddeg}. The paper does not say
whether its extension-range counts use only the classes with determined optimality. For $E[\ell]$ irreducible the question does not affect
$v_\ell$ of the degree, by (iv), nor the terms $e_\ell$, $t_\ell$, $h_\ell$; but the weight $w_3$ is only observed, not proved, to be the same for all
curves of a class (\Cref{ch:ellcurves}).
\item \texttt{ellweilcurve} identifies the optimal curve inside PARI for moderate conductors (\Cref{rem:computation:ellmoddegree}(ii)), and
\texttt{opt\_man} records it (\Cref{ex:computation:curvedata}).
\item When $E[\ell]$ is irreducible, every isogeny in the class has degree prime to $\ell$, so $v_\ell$ of the degree of the minimal
parametrisation is the same for every curve of the class; this is the paper's remark in Section 1 and is proved in \Cref{ch:moddeg}. When
$E[\ell]$ is reducible it is not: the minimal parametrisations of \curve{11a1}, \curve{11a2}, \curve{11a3} have degrees $1$, $5$, $5$ (\texttt{alldegphi}), and
\texttt{rule}, which uses $\deg\varphi/c_\varphi^2$, returns \texttt{[0, 1]}, \texttt{[1, 0]}, \texttt{[-1, 0]} at $\ell=5$.
\end{enumerate}
\end{remark}

\begin{convention}[labels]\label{conv:computation:labels}
\begin{enumerate}[label=(\roman*)]
\item A Cremona label is the conductor, a class code (\texttt{a}, \dots, \texttt{z}, \texttt{ba}, \texttt{bb}, \dots) and the number of the curve in its
class: \curve{11a1}, \curve{539c1}; \curve{121c} and \curve{121d} denote classes. The capitalised form \curve{99A1} of \cite{Cremona97} and \cite{ARS12} is
\curve{99a1}.
\item An LMFDB label \cite{LMFDB} is $N$.class number, with classes and curves ordered differently; the file \texttt{alllabels} is the
dictionary. For example \curve{11a1} is \texttt{11.a2} and \curve{726a1} is \texttt{726.b2}.
\item The images in \texttt{galrep} carry Sutherland's labels \cite{Sutherland16}, whose general scheme and table at $\ell=3$ are in
\Cref{ch:galrep}. What this chapter needs: an image is reducible if and only if its label begins with ``$\ell$B'' (contained in a Borel
subgroup) or ``$\ell$Cs'' (contained in a split Cartan subgroup), and the prime part of a label must be compared exactly (the label
\texttt{37B.8.1} is a $37$-adic label and must not be read as a $3$-adic one). The three labels at $\ell=3$ that the paper names are \texttt{3Ns}
(image $N_s(\F_3)$), \texttt{3Cs.1.1} ($E[3]\simeq\Z/3\oplus\mu_3$, image the subgroup of order $2$ of the split Cartan subgroup that fixes a point) and
\texttt{3B.1.1} (a rational point of order $3$, image in a Borel subgroup, $E[3]$ not split). Among the $37{,}853$ optimal curves without CM below
$10{,}000$ the $3$-adic labels occur as follows: \texttt{3B.1.1} $1{,}858$, \texttt{3B} $1{,}567$, \texttt{3B.1.2} $503$, \texttt{3Nn} $170$, \texttt{3Ns} $90$,
\texttt{3Cs.1.1} $49$, \texttt{3Cs} $25$. These are counts of optimal curves; counts over all curves of the tables, as in \Cref{ch:galrep}, are larger.
\item \texttt{galrep} omits the images forced by complex multiplication (\Cref{ex:computation:mod3-image}); CM curves are therefore recognised
by $j$, never by the absence of labels.
\end{enumerate}
\end{convention}

\begin{paperlink}\label{pl:computation:ecdata}
Section sec:data: ``We used Cremona's tables \cite{Cremona97} in the form of the \texttt{ecdata} repository at commit \texttt{25cec5e}''. The six
files used here are byte for byte those of that commit (their SHA-256 sums are in \Cref{conv:computation:checklist}), so the book's counts
below conductor $10{,}000$ are computed from the paper's own data. Section 1.4 reads the component groups off the Kodaira types; the paper
recomputes the types with PARI, and \Cref{ex:computation:curvedata}(a) shows that, below conductor $10{,}000$, the repository's own
\texttt{curvedata} agrees.
\end{paperlink}
\section{The paper's function \texttt{rule(E, l, iii)}, line by line}\label{sec:computation:rule}

The paper's appendix prints two GP functions, \texttt{rule} and \texttt{ns3}, which check one curve against (eq:rule). This section prints them
verbatim, reads them line by line as code, proves that they compute the right side of (eq:rule) and of its III-free form under stated
conditions, and lists the cases in which they return something other than the paper's quantity. The mathematics of each term is in
\Cref{ch:lfunctions,ch:paper}; here we read the program.

\begin{construction}[the appendix code]\label{con:computation:rule-listing}
The functions of the paper's appendix, which the appendix says ``are the file \texttt{scripts/euler\_rule.gp}'', copied verbatim, with line
numbers added:
\begin{lstlisting}[language=GP,numbers=left,numberstyle=\tiny,basicstyle=\ttfamily\scriptsize]
rule(E, l, {iii = 0}) =
{ my(R = ellglobalred(E), P = R[4][,1], K = R[5], S = lfunsympow(E, 2), b = 0, j = E.j);
  for(i = 1, #P,
    my(q = P[i], k = K[i][2], e = 0, t = 0, vj = valuation(j, q));
    if(q != l,
      my(F = 1/lfuneuler(S, q), r = subst(F, x, 1/q^2)); \\ P_q(q^{-2}): omitted factor
      if(k >= 5, r /= (1 - 1/q^2));                     \\ multiplicative: naive one out
      if(abs(k) != 3 || iii, e = valuation(r, l)),      \\ III, III* (codes 3,-3) if iii
      if(l == 3, e = if(k == 4 || k == 2, 1,
                     if(k == -4 || k == -2, 2, if(k == -3, 1, 0)))));   \\ w_3 at q = 3
    if(vj < 0, t = valuation(-vj, l));                   \\ Tamagawa exponent
    b += e + t);
  if(iii && l == 3, b -= ns3(E));                        \\ the correction term h_3
  [valuation(ellmoddegree(E), l), b]; }
ns3(E) = my(F = factor(elldivpol(E, 3))); (#F[,1] == 2 && poldegree(F[1,1]) == 2);
\end{lstlisting}
The fifteen lines contain no double hyphen. They are loaded with \texttt{read("euler\_rule.gp")} in \texttt{gp}, or with
\texttt{pari('read("euler\_rule.gp")')} in Python (\Cref{warn:computation:pitfalls}(g)). The appendix adds: ``Here \texttt{K[i][2]} is
PARI's Kodaira code: $n+4$ for $\mathrm{I}_n$ and $-(n+4)$ for $\mathrm{I}_n^*$ when $n\ge1$, $1$ for $\mathrm{I}_0$ and $-1$ for $\mathrm{I}_0^*$, $\pm2$ for
$\mathrm{II}/\mathrm{II}^*$, $\pm3$ for $\mathrm{III}/\mathrm{III}^*$, $\pm4$ for $\mathrm{IV}/\mathrm{IV}^*$. The function \texttt{rule(E,l)} leaves the type
$\mathrm{III}$ terms out and \texttt{rule(E,l,1)} is (eq:rule)''.
\end{construction}

\begin{remark}[reading the code]\label{rem:computation:rule-reading}
\begin{enumerate}[label=(\alph*)]
\item Lines 1 and 2. \texttt{rule} takes a curve, an odd prime and an optional flag \texttt{iii}, $0$ by default. \texttt{R} is the output of
\texttt{ellglobalred} (\Cref{def:computation:ellglobalred}); \texttt{P = R[4][,1]} is the column of bad primes in increasing order, and
\texttt{K = R[5]} the local data in the same order, so that \texttt{K[i]} refers to \texttt{P[i]}. \texttt{S} is the symmetric square
(\Cref{rem:computation:lfunsympow}), \texttt{b} the bound being accumulated, \texttt{j} the $j$-invariant.
\item Lines 3 and 4. A loop over the bad primes; for each, \texttt{k} is the Kodaira code, \texttt{e} and \texttt{t} the two terms (both $0$ unless
set), and \texttt{vj} $=v_q(j)$, which is \texttt{+oo} when $j=0$.
\item Line 6, for $q\ne\ell$. \texttt{F = 1/lfuneuler(S, q)} is the polynomial $P_q(X)$ in the variable \texttt{x}, and \texttt{r} the exact
rational $P_q(q^{-2})=L_q(A_f,1)^{-1}$.
\item Line 7. The test \texttt{k >= 5} selects the multiplicative primes (\Cref{conv:computation:code-idioms}). There the naive factor at $s=1$ is
$1-a_q^2q^{-2}=1-q^{-2}$, since $a_q=\pm1$, and the division makes \texttt{r} the quotient of Definition def:terms; it equals $1$, because
$P_q(X)=1-X$ at a multiplicative prime, so $e_\ell(E,q)=0$ there. At an additive prime $a_q=0$, the naive factor is $1$, and no division is
needed.
\item Line 8. The valuation of \texttt{r} is taken unless the type is $\kod{III}$ or $\kod{III}^*$ (\texttt{abs(k) == 3}) and \texttt{iii} is $0$: so
\texttt{rule(E, l)} leaves the type $\kod{III}$ terms at $q\ne\ell$ out, and \texttt{rule(E, l, 1)} keeps them.
\item Lines 9 and 10, the \texttt{else} branch of line 5, for $q=\ell$. No Euler term is taken. For $\ell=3$ the variable \texttt{e} is set to the
weight $w_3$ read off the code: $1$ for $\kod{IV}$ and $\kod{II}$ (codes $4$, $2$), $2$ for $\kod{IV}^*$ and $\kod{II}^*$ ($-4$, $-2$), $1$ for $\kod{III}^*$ ($-3$), and $0$ for
every other type ($\kod{III}$, $\kod{I}_0^*$, and the potentially multiplicative types, whose contribution is $t_3$). For $\ell\ge5$ nothing is added
at $q=\ell$. The flag \texttt{iii} does not act here: \texttt{rule(E, 3)} counts the weight $1$ of a type $\kod{III}^*$ at $q=3$.
\item Line 11. For every bad prime, $q=\ell$ included, \texttt{t} $=v_\ell(-v_q(j))$ if $v_q(j)<0$; this is $t_\ell(E,q)$ of Definition def:terms, which
the paper defines at $q=\ell$ too. When $j=0$, \texttt{vj} is \texttt{+oo}, the test fails and \texttt{t} $=0$, as it should: $j=0$ means potentially good
reduction everywhere.
\item Line 12 adds the two terms; line 13 subtracts \texttt{ns3(E)} when \texttt{iii} is set and $\ell=3$; for $\ell\ge5$ nothing is subtracted, as
$h_\ell(E)=0$ for $\ell\ge5$ (\Cref{ch:galrep}).
\item Line 14 returns the pair $[v_\ell(\texttt{ellmoddegree(E)}),\ b]$. \texttt{ellmoddegree} is by far the slowest step
(\Cref{rem:computation:ellmoddegree}(iv)).
\item Line 15. \texttt{ns3(E)} factors $\psi_3$ and returns $1$ if the factorisation has exactly two distinct irreducible factors, the first of
degree $2$, and $0$ otherwise.
\end{enumerate}
\end{remark}

\begin{proposition}[what \texttt{rule} computes]\label{prop:computation:rule-correct}
Let $E/\Q$ be an elliptic curve given by any Weierstrass model, and $\ell$ an odd prime. Assume that \texttt{lfuneuler(lfunsympow(E, 2), q)}
returns $1/P_q(X)$ and that \texttt{ellmoddegree(E)} returns $\deg\varphi/c_\varphi^2$ for the minimal parametrisation $\varphi$ of the given curve, as
PARI's documentation states.
\begin{enumerate}[label=(\roman*)]
\item For every bad prime $q\ne\ell$, the number \texttt{e} that \texttt{rule(E, l, 1)} computes at $q$ is $e_\ell(E,q)$; at every bad prime, the
number \texttt{t} is $t_\ell(E,q)$; at $q=\ell=3$, \texttt{e} is $w_3(E)$.
\item \texttt{ns3(E)} $=1$ if and only if $\psi_3$ is a product of two irreducible quadratic polynomials over $\Q$. If $E[3]$ is irreducible, then
\texttt{ns3(E)} $=h_3(E)$, and both are $1$ exactly when the mod $3$ image is $N_s(\F_3)$.
\item Hence, if $E$ is the optimal curve of its class (with Manin constant $1$) and, when $\ell=3$, $E[3]$ is irreducible, then
\texttt{rule(E, l, 1)} returns $[v_\ell(\degphi),B^{\mathrm{rule}}_\ell(E)]$ and \texttt{rule(E, l)} returns $[v_\ell(\degphi),B^{\mathrm{III}}_\ell(E)]$; and
$B^{\mathrm{nc}}_\ell(E)$ is the second entry of \texttt{rule(E, l, 1)} plus \texttt{ns3(E)} at $\ell=3$, and equal to it at $\ell\ge5$.
\end{enumerate}
\end{proposition}

\begin{proof}
(i) The local data are those of a minimal model whatever model is given (\Cref{def:computation:ellglobalred}), and the code \texttt{k} is $\ge5$
exactly at the multiplicative primes, the code $1$ of good reduction never occurring at a bad prime. At a multiplicative prime $a_q^2=1$ and the
naive factor at $s=1$ is $1-q^{-2}$; at an additive prime $a_q=0$ and the naive factor is $1$ (\Cref{ch:lfunctions}). So, by
\Cref{rem:computation:rule-reading}(c), (d), after line 7 \texttt{r} is $P_q(q^{-2})/(1-a_q^2q^{-2})$, and line 8 with \texttt{iii} $=1$ takes its
$\ell$-adic valuation: this is the definition of $e_\ell(E,q)$. Line 11 is the definition of $t_\ell(E,q)$ word for word, the case $j=0$ being
handled by \texttt{valuation(0, q)} $=$ \texttt{+oo}. Lines 9 and 10 are the definition of $w_3(E)$ in Conjecture conj:euler: the weight is defined
for potentially good reduction at $3$, and the code gives $0$ to the potentially multiplicative types, which are not potentially good.

(ii) PARI factors an integer polynomial into irreducible primitive factors and drops the content (\Cref{warn:computation:pitfalls}(f)), so the
leading coefficient $3$ of $\psi_3$ produces no extra factor. Since $\psi_3$ has distinct roots, its factorisation has no repeated factor. If
there are exactly two factors and the first has degree $2$, the other has degree $4-2=2$; conversely, a factorisation $2+2$ has exactly two
factors and both have degree $2$, whatever order PARI lists them in. So \texttt{ns3(E)} $=1$ if and only if the pattern is $2+2$. If $E[3]$ is
irreducible, $\psi_3$ has no rational root (a rational root gives a Galois-stable subgroup of order $3$, the $x$-coordinate determining the
point up to sign), so the possible patterns are $4$ and $2+2$, and the pattern is $2+2$ exactly when the image is $N_s(\F_3)$ (Proposition prop:ns;
\Cref{ch:galrep}). For $E[3]$ irreducible, $h_3(E)\le1$, with equality exactly when $\bar\rho\simeq\bar\rho\otimes\chi_{-3}$, which happens exactly for image
$N_s(\F_3)$, with or without CM (\Cref{ch:galrep}, where the strengthened Lemma lem:h0 and Proposition prop:ns are proved).

(iii) By (i) the sum accumulated in \texttt{b} by line 12 is $B^{\mathrm{nc}}_\ell(E)$ when \texttt{iii} $=1$ and $B^{\mathrm{III}}_\ell(E)$ when \texttt{iii} $=0$. Line 13
subtracts \texttt{ns3(E)} $=h_3(E)$ in the first case at $\ell=3$, by (ii), giving $B^{\mathrm{rule}}_3(E)$; at $\ell\ge5$, $h_\ell(E)=0$ (\Cref{ch:galrep}), so
$B^{\mathrm{rule}}_\ell=B^{\mathrm{nc}}_\ell$. Finally, for the optimal curve with $c_E=1$, \texttt{ellmoddegree(E)} $=\degphi$, so the first entry is $v_\ell(\degphi)$.
\end{proof}

The proposition is about PARI's functions as documented; the rest of this chapter checks them against independent data. Below conductor $10{,}000$
the book's Python implementation of the five bounds (\Cref{rem:computation:book-scripts}), which reads the stored factors and codes and never
calls \texttt{rule}, agrees with the pure GP loop of \Cref{con:computation:table-small} on every optimal curve of conductor below $2{,}000$ at
$\ell=3,5,7$ (\Cref{ex:computation:repro-uniform}), and \texttt{ellmoddegree} agrees with Cremona's degrees on all $38{,}042$ optimal curves
(\Cref{ex:computation:moddeg-check}).

\begin{warning}[where \texttt{rule} does not compute the paper's quantity]\label{warn:computation:rule-pitfalls}
\begin{enumerate}[label=(\alph*)]
\item \emph{A non-optimal curve.} \texttt{ellmoddegree} returns $\deg\varphi/c_\varphi^2$ for the curve given, so the first entry is not $v_\ell(\degphi)$:
on \curve{11a3} at $\ell=5$, \texttt{rule} returns \texttt{[-1, 0]} (a negative ``valuation of a degree''), and on \curve{990h1} at $\ell=3$ it returns
\texttt{[4, 3]} where the optimal \curve{990h3} returns \texttt{[3, 3]}. Curve number $1$ is not always the optimal curve
(\Cref{rem:computation:optimality-data}).
\item \emph{CM curves.} \texttt{ns3} is $1$ on $13$ optimal CM curves below $10{,}000$ (\Cref{ex:computation:mod3-image}). Conjecture conj:euler excludes
CM, so these outputs concern no statement of the paper.
\item \emph{Eisenstein curves.} \curve{11a1} at $\ell=5$ returns \texttt{[0, 1]}, an apparent violation; but $E[5]$ is reducible (\texttt{5Cs.1.1}) and the
curve is outside every statement of the paper (\Cref{ex:computation:11a1-eisenstein}). The population filter is part of the statement.
\item \emph{$\ell\mid N$.} \texttt{rule} runs, and the data satisfy (eq:rule) there, but no theorem applies: Theorem thm:main needs $\ell\nmid N$.
\item \emph{Codes.} $n+4$ and $-(n+4)$ are the codes of $\kod{I}_n$ and $\kod{I}_n^*$ only for $n\ge1$ (\Cref{conv:computation:code-idioms}).
\item \emph{Versions.} In 2.17.2 \texttt{ellmoddegree} returns an integer on an optimal curve and a fraction on some non-optimal ones
(type \texttt{t\_INT} and \texttt{t\_FRAC}); \texttt{valuation} accepts both. The appendix says the functions run with ``version 2.15 or later''; we
have not run 2.15. The book's own scripts also use \texttt{readstr} and \texttt{strsplit}, which exist in 2.17.2 and whose presence in 2.15 we
have not checked.
\item \emph{Speed.} For a table, take $\degphi$ from \texttt{degphi} and compute only the bound: the function \texttt{rulebound} of
\Cref{con:computation:table-small}.
\end{enumerate}
\end{warning}

\begin{paperlink}\label{pl:computation:rule}
The appendix: ``A single curve can be checked against (eq:rule) in PARI/GP (version 2.15 or later) with the following functions, which
return $[v_\ell(\degphi),\text{bound}]$ \dots\ The curve must be the optimal curve of its class and $E[\ell]$ must be irreducible: for instance
990h1 returns \texttt{[4, 3]} while the optimal curve 990h3 of its class returns \texttt{[3, 3]}, and 11a1 at the Eisenstein prime $\ell=5$ returns
\texttt{[0, 1]}.'' \Cref{prop:computation:rule-correct} is the precise form of this, and \Cref{warn:computation:rule-pitfalls}(a) and (c) reproduce
the two examples. Optimality is what makes the first entry $v_\ell(\degphi)$; irreducibility of $E[3]$ is what makes \texttt{ns3} the correction at
$\ell=3$; at $\ell\ge5$ irreducibility is needed not for the computation but for any statement of the paper to apply.
\end{paperlink}

\section{Worked sessions}\label{sec:computation:sessions}

Each session below gives the input, the exact output of PARI/GP 2.17.2, the per-prime breakdown printed by \texttt{rule\_terms(E, l, iii)}
(\Cref{con:computation:sessions}), whose rows are $[q,\allowbreak k,\allowbreak P_q(X),\allowbreak P_q(q^{-2}),\allowbreak e,w,t]$ (the Euler term $e$ for $q\ne\ell$, the weight $w$ at $q=\ell=3$,
the term $t$; \texttt{iii} defaults to $1$), and a short reading. The mathematics of the terms is in \Cref{ch:lfunctions,ch:paper}. In the paper
these are the examples of the appendix, the Example of Section sec:proof, the Remark after it, Remark rem:wild and Section sec:tam.

\begin{example}[\curve{121a1}]\label{ex:computation:121a1}
\begin{lstlisting}[language=GP]
? read("euler_rule.gp"); read("sessions.gp");
? E = ellinit([1,1,1,-30,-76]);
? [rule(E, 3), rule(E, 3, 1), ellmoddegree(E)]
[[1, 1], [1, 1], 6]
? rule_terms(E, 3)
[[11, 2, 11*x + 1, 12/11, 1, 0, 0]]
\end{lstlisting}
The only bad prime is $11$, of type $\kod{II}$ (code $2$), with $P_{11}(X)=1+11X$ and $e_3(E,11)=v_3(12/11)=1$; $t_3=0$ since $v_{11}(j)=1\ge0$. So the bound is
$1$, and $\degphi=6=2\cdot3$ attains it. The paper's appendix: ``For example \texttt{rule(ellinit([1,1,1,-30,-76]), 3)} (the curve 121a1)
returns \texttt{[1, 1]}''. Here $B^{\mathrm{main}}_3=0$ ($\#\Phig_{11}=1$ for type $\kod{II}$) while $B^{\mathrm{tw}}_3=1$: this curve is the simplest case where (eq:twist)
is stronger than (eq:main).
\end{example}

\begin{example}[\curve{726a1}, the correction term]\label{ex:computation:726a1}
\begin{lstlisting}[language=GP]
? E = ellinit([1,1,0,-35,-51]);
? [rule(E, 3), rule(E, 3, 1), ns3(E), ellmoddegree(E)]
[[2, 2], [2, 2], 1, 144]
? rule_terms(E, 3)
[[2, 10, -x + 1, 3/4, 0, 0, 1], [3, 7, -x + 1, 8/9, 0, 0, 1], [11, 3, 11*x + 1, 12/11, 1, 0, 0]]
? rule_terms(E, 3, 0)
[[2, 10, -x + 1, 3/4, 0, 0, 1], [3, 7, -x + 1, 8/9, 0, 0, 1], [11, 3, 11*x + 1, 12/11, 0, 0, 0]]
\end{lstlisting}
At $2$, type $\kod{I}_6$: $e=0$, $t_3=v_3(6)=1$. At $3=\ell$, type $\kod{I}_3$: no Euler term, $w=0$, $t_3(E,3)=v_3(3)=1$. At $11$, type $\kod{III}$, $P_{11}=1+11X$,
$e_3=v_3(12/11)=1$, kept only with \texttt{iii}. So $B^{\mathrm{III}}_3=2$, $B^{\mathrm{nc}}_3=3$, \texttt{ns3} $=1$ ($\psi_3=(x^2-x-74)(3x^2+8x+20)$,
\Cref{ex:computation:mod3-image}) and $B^{\mathrm{rule}}_3=2$; and $\degphi=144=2^43^2$. The appendix: ``for 726a1, a curve with image $\Ns$ and type
$\mathrm{III}$ at $11$, \texttt{rule(E,3)} and \texttt{rule(E,3,1)} both return \texttt{[2, 2]}, the type $\mathrm{III}$ term and the correction
cancelling.'' Without the correction the bound $B^{\mathrm{nc}}_3=3$ fails by exactly one, and the failure involves the prime $2\equiv-1\pmod3$ with
$t_3(E,2)=1$, as Section sec:anomaly says of every such failure. The curve has $3\mid N$, so it is outside Theorem thm:main in any case; and if
(eq:thm) is read literally, with terms at $q\ne3$ only, its bound is $B^{\mathrm{thm}}_3=1+1=2$ and there is no violation
(\Cref{ex:computation:repro-ns}). The congruence number $r_E=1584=11\cdot144$ (\Cref{ch:moddeg}) is not computed by \texttt{rule}.
\end{example}

\begin{example}[\curve{539c1}]\label{ex:computation:539c1}
\begin{lstlisting}[language=GP]
? E = ellinit([1,0,1,170,-3237]);
? [rule(E, 3), rule(E, 5), ellmoddegree(E)]
[[2, 2], [0, 0], 288]
? rule_terms(E, 3)
[[7, -7, -x + 1, 48/49, 1, 0, 1], [11, 6, -x + 1, 120/121, 0, 0, 0]]
\end{lstlisting}
At $7$, code $-7$, type $\kod{I}_3^*$: $P_7=1-X$, $P_7(7^{-2})=48/49$, no division (the prime is additive), $e_3=v_3(48)=1$, and $t_3=v_3(3)=1$. At $11$, type $\kod{I}_2$:
nothing. So the bound is $2=v_3(288)$. The Tate parameter at $7$ is that of the twist \curve{77c1}, with $\bar u=5$ not a cube in $\F_7^\times$
(\Cref{ex:computation:tate-parameter}), so the Tamagawa factor of the adjoint adds nothing beyond $t_3$ (Lemma lem:tamfactor). The paper (Example of
Section sec:proof): ``$\deg\varphi_{\text{539c1}}=2^5\cdot3^2$, and the two terms $e_3(E,7)=v_3(48)=1$ and $t_3(E,7)=v_3(3)=1$ are both attained.''
\end{example}

\begin{example}[\curve{980i1}]\label{ex:computation:980i1}
\begin{lstlisting}[language=GP]
? E = ellinit([0,0,0,1568,-72716]);
? [rule(E, 3), rule(E, 5), ellmoddegree(E)]
[[2, 2], [1, 1], 2880]
? rule_terms(E, 3)
[[2, -4, 2*x + 1, 3/2, 1, 0, 0], [5, 5, -x + 1, 24/25, 0, 0, 0], [7, -9, -x + 1, 48/49, 1, 0, 0]]
? rule_terms(E, 5)
[[2, -4, 2*x + 1, 3/2, 0, 0, 0], [5, 5, -x + 1, 24/25, 0, 0, 0], [7, -9, -x + 1, 48/49, 0, 0, 1]]
\end{lstlisting}
At $\ell=3$: type $\kod{IV}^*$ at $2$ with $P_2=1+2X$ and $e_3=v_3(3/2)=1$; type $\kod{I}_5^*$ at $7$ with $e_3=v_3(48/49)=1$. At $\ell=5$: $t_5(E,7)=v_5(5)=1$, all $e_5=0$.
$\degphi=2880=2^6\cdot3^2\cdot5$, and both bounds are attained. The contribution of $2$ at $\ell=3$ is a theorem, not an observation: a type $\kod{IV}$ or
$\kod{IV}^*$ at $2$ always has $P_2=1+2X$ (Lemma lem:IV; \Cref{ch:lfunctions}), and the check of \Cref{ex:computation:repro-further}(c) confirms it on every
optimal curve below $10{,}000$. At $\ell=5$ the prime $5$ divides $N$, so the curve illustrates the congruence of the Example of Section sec:proof, not
Theorem thm:qual (as the paper says).
\end{example}

\begin{example}[\curve{99a1}]\label{ex:computation:99a1}
\begin{lstlisting}[language=GP]
? E = ellinit([1,-1,1,-2,0]);
? [rule(E, 3), ellmoddegree(E)]
[[0, 0], 4]
? rule_terms(E, 3)
[[3, 3, 3*x + 1, 4/3, 0, 0, 0], [11, 5, -x + 1, 120/121, 0, 0, 0]]
\end{lstlisting}
$N=3^2\cdot11$. At $3=\ell$, type $\kod{III}$ (code $3$): no Euler term, $w_3=0$, and $t_3=0$ since $v_3(j)=9\ge0$; PARI's $P_3=1+3X$ is printed by
\texttt{rule\_terms} but not used. At $11$, type $\kod{I}_1$. The bound is $0$ and $\degphi=4$. This is the curve \curve{99A1} of \cite{ARS12}, whose congruence number is
$r_E=12$ (\Cref{ch:moddeg}): the gap at $3$ between $r_E$ and $\degphi$ is allowed by the theorem of Agashe, Ribet and Stein because $9\mid N$
(\Cref{rem:computation:congruence-number}).
\end{example}

\begin{example}[the prime $2$ through PARI's factor]\label{ex:computation:wild-two}
\begin{lstlisting}[language=GP]
? E = ellinit([0,1,0,-1,2]);
? [rule(E, 3), ellmoddegree(E)]
[[1, 1], 24]
? rule_terms(E, 3)[1]
[2, 2, 2*x + 1, 3/2, 1, 0, 0]
? E = ellinit([0,0,0,32,-16]);
? [rule(E, 3), ellmoddegree(E)]
[[2, 2], 288]
? rule_terms(E, 3)[1]
[2, -2, -8*x^3 - 4*x^2 + 2*x + 1, 9/8, 2, 0, 0]
? factor(-8*x^3 - 4*x^2 + 2*x + 1)
[2*x - 1, 1; 2*x + 1, 2]
? E = ellinit([0,0,0,-16,-16]);
? [rule(E, 5), ellmoddegree(E)]
[[1, 1], 80]
? rule_terms(E, 5)
[[2, -2, -8*x^3 + 4*x^2 - 2*x + 1, 5/8, 1, 0, 0], [37, 5, -x + 1, 1368/1369, 0, 0, 0]]
? factor(-8*x^3 + 4*x^2 - 2*x + 1)
[2*x - 1, 1; 4*x^2 + 1, 1]
\end{lstlisting}
\curve{336f1} $=[0,1,0,-1,2]$ has type $\kod{II}$ at $2$ and nevertheless $P_2=1+2X$, $e_3(E,2)=1$, bound $1$, $\degphi=24$. \curve{1232l1} $=[0,0,0,32,-16]$ has type $\kod{II}^*$ at
$2$ and the cubic $P_2=(1-2X)(1+2X)^2$, $e_3(E,2)=v_3(9/8)=2$, bound $2$, $\degphi=288$: the prime $2$ alone attains the bound. \curve{592c1} $=[0,0,0,-16,-16]$ has
type $\kod{II}^*$ at $2$, $P_2=(1-2X)(1+4X^2)$, $e_5(E,2)=v_5(5/8)=1$, bound $1$ at $\ell=5$, $\degphi=80=2^4\cdot5$: the prime $2$ contributes at $\ell=5$. At $2$ the
Kodaira type does not determine $P_2$; the four possible shapes are proved in \Cref{ch:lfunctions} (Proposition prop:shapes). Among the optimal
curves without CM below conductor $10{,}000$, type $\kod{II}$ at $2$ gives $P_2=1$ for $1{,}092$ curves, $1+2X$ for $371$ and a cubic for $74$; type $\kod{II}^*$
gives $514$, $74$ and $191$. \emph{A data correction.} Remark rem:wild says: ``Among optimal non-CM curves of conductor below $10{,}000$, every cubic factor
at $2$ ($1{,}143$) and at $3$ ($2{,}283$) is of shape (ii).'' The two counts are over all $64{,}687$ curves of the tables, not over the optimal curves
without CM, among which there are $534$ cubic factors at $2$ and $1{,}308$ at $3$ (recomputed here with \texttt{lfunsympow}; \Cref{ch:lfunctions} found the
same). Likewise the remark's ``$58\%$'' of potentially good types $\kod{I}_n^*$ at $2$ is the share among optimal curves ($2{,}340$ of $4{,}037$); over all
curves it is $61\%$ ($5{,}935$ of $9{,}667$). And ``of shape (ii)'' is automatic, a cubic $P_q$ having no other shape; the empirical content is that each
cubic equals the formula of (ii) for a twist of good reduction found independently, which \Cref{ch:lfunctions} checks.
\end{example}

\begin{example}[\curve{11a1} at an Eisenstein prime]\label{ex:computation:11a1-eisenstein}
\begin{lstlisting}[language=GP]
? E = ellinit([0,-1,1,-10,-20]);
? [rule(E, 5), rule_terms(E, 5)]
[[0, 1], [[11, 9, -x + 1, 120/121, 0, 0, 1]]]
\end{lstlisting}
The bound is $t_5(E,11)=v_5(5)=1$ and $\degphi=1$: an apparent violation. But \texttt{galrep} gives \texttt{5Cs.1.1}, so $E[5]\simeq\Z/5\oplus\mu_5$ is reducible, and
the curve belongs to no statement of the paper except Conjecture conj:eis. There the deficit is $D=v_5(\#\Phig_{11})-v_5(\degphi)=1-0=1\le v_5(\#E(\Q)_{\mathrm{tors}})=1$.
The Tate data (\Cref{ex:computation:tate-parameter}) say more: $\bar u=10\equiv-1$ is a fifth power in $\F_{11}^\times$, so the formula of Lemma lem:tamfactor would give
$v_5(\Tam_{11}(\ad T))=1+\min(1,1,1)=2$, while $v_5(\degphi)=0$. The identity (eq:identity) fails here because its hypothesis fails: the population filter is
part of every statement, not an afterthought (\Cref{ch:cohomology,ch:blochkato} make the same point in cohomology).
\end{example}

\begin{example}[the Tamagawa refinement on three curves]\label{ex:computation:tamagawa-session}
At $\ell=3$, for a multiplicative prime $q\equiv1\pmod3$ with $3\mid n$, Lemma lem:tamfactor adds $\min(v_3(n),v_3(q-1),v_3(c(u)))$ to the Tamagawa
exponent (Section sec:tam):
\begin{lstlisting}[language=GP]
? E = ellinit([1,0,1,-31,117]);
? [rule(E, 3), ellmoddegree(E), ellglobalred(E)[4][,1]~]
[[2, 1], 144, [7, 113]]
? E = ellinit([1,1,0,4,11]);
? [rule(E, 3), ellmoddegree(E), ellglobalred(E)[4][,1]~]
[[1, 1], 6, [7, 11]]
? E = ellinit([0,-1,1,30,-335]);
? [rule(E, 3), ellmoddegree(E), ellglobalred(E)[4][,1]~]
[[2, 1], 720, [19, 83]]
\end{lstlisting}
\curve{791a1} at $7$: $n=3$, $\bar u=6$ is a cube (\Cref{ex:computation:tate-parameter}), the extra term is $\min(1,1,1)=1$, so the refined bound is $1+1=2=v_3(144)$:
attained. \curve{77c1} at $7$: $\bar u=5$ is not a cube, the extra term is $0$, and the bound $1$ is attained by $\degphi=6$. \curve{1577a1} at $19$: $\bar u=8$ is a
cube but not a ninth power, $v_3(18)=2$, the extra term is $\min(1,2,1)=1$, and the refined bound $2=v_3(720)$ is attained. The cohomology is in
\Cref{ch:galcoh,ch:blochkato}.
\end{example}
\section{Computing modular degrees and congruence numbers}\label{sec:computation:moddeg}

This short section says how the modular degree is computed, why it is the expensive step, how the book checked Cremona's degrees, and where
the congruence number, which PARI does not compute, comes from. In the paper it is the sentence ``The modular degree is computable
\cite{Cremona97,Watkins02}'' of Section 1, the recomputations of Section sec:data and Section sec:oot, and the Remark after the Example of Section
sec:proof.

\begin{remark}[three methods]\label{rem:computation:moddeg-methods}
\emph{Zagier's method} \cite{Zagier85} computes the map $\varphi_{E*}$ on $H_1(X_0(N),\Z)$ with modular symbols and reads $\degphi$ off an index in
homology; \emph{Cremona's} \cite[Chapter II]{Cremona97} computes the period lattice of $E$ from modular symbols and compares it with the lattice
of $\varphi_E^*$, which also certifies optimality. \emph{Watkins's method} \cite{Watkins02} uses the formula $\degphi\cdot\vol(\Lambda_E)=4\pi^2c_E^2\langle f,f\rangle$
together with the expression of $\langle f,f\rangle$ through $L(\Sym f,2)$ and the bad Euler factors; it is the method of \texttt{ellmoddegree}. The
mathematics of all three is in \Cref{ch:moddeg,ch:lfunctions}. Why the computation is expensive (a heuristic, not a statement about PARI's
code): $L(\Sym f,2)$ is evaluated by Dokchitser's method \cite{Dokchitser04}, which needs a number of Dirichlet coefficients roughly proportional
to the square root of the conductor of $L(\Sym f,s)$; that conductor is at most $N^2$ when $(N,6)=1$ (\Cref{ch:lfunctions}), so the work grows
roughly linearly with $N$, and the coefficients themselves cost more as $N$ grows. Cremona's tables contain the degrees of all optimal curves
up to conductor $500000$; the paper recomputed fifteen of them inside the tables and $893$ outside (Section sec:oot). The congruence module of
\Cref{ch:moddeg} is a Smith normal form computation (\texttt{matsnf}, \texttt{mathnf} in PARI; \Cref{ch:modules}), and the modular-forms
functions \texttt{mfinit}, \texttt{mfbasis}, \texttt{mfcoefs}, \texttt{mfheckemat}, \texttt{mfsturm}, \texttt{mfpetersson} used in \Cref{ch:modforms,ch:hecke,ch:moddeg,ch:modgalrep}
are explained there with their normalisations. One trap is recorded in \Cref{warn:computation:pitfalls}(h): \texttt{mfsturm} applied to a space returns less than Sturm's bound.
\end{remark}

\begin{example}[\texttt{ellmoddegree} against Cremona's degrees]\label{ex:computation:moddeg-check}
The script \texttt{moddeg\_check.py} (\Cref{rem:computation:book-scripts}) runs \texttt{ellmoddegree} on the equation of each of the $38{,}042$
lines of \texttt{degphi} and compares with the degree written there, in two processes, with \texttt{parisizemax} $2\cdot10^9$, catching errors.
MODDEGRESULT
The paper's sentence ``fifteen modular degrees spanning the configurations were recomputed with \texttt{ellmoddegree}, all agreeing with the
tables'' (Section sec:data) is thereby extended, in the book's range, to every optimal curve; the paper's revised version says the same of its
own check below $10{,}000$.
\end{example}

\begin{remark}[the congruence number]\label{rem:computation:congruence-number}
PARI has no function for the congruence number $r_E$. \Cref{ch:moddeg} computes it from \texttt{mfinit}, \texttt{mfbasis} and the Hecke action on
the lattice of integral forms, and finds $r_E=12$ for \curve{99a1} against $\degphi=4$, $r_E=1584=11\cdot144$ for \curve{726a1}, and $r_E=504=7\cdot72$ for
\curve{539d1}. The theorem of Agashe, Ribet and Stein \cite[Theorem 2.2]{ARS12} gives $\degphi\mid r_E$, with equal $p$-parts when $p^2\nmid N$; in the
three examples the extra prime ($3$, $11$, $7$) has square dividing $N$. This is the Remark after the Example of Section sec:proof (``Table 1 of
\cite{ARS12} lists the curve 99A1, of conductor $9\cdot11$, with modular degree $4$ and congruence number $12$'').
\end{remark}

\section{Populations, violations and descriptive statistics}\label{sec:computation:populations}

Before any test, one must say what is counted, in which population, and what a rate is a rate of. This section fixes the units and the
descriptive statistics of the paper's tables (equality rates, mean excess, base rates) and computes them on the book's range. In the paper it
is the background of every table of Section sec:data.

\begin{definition}[units and populations]\label{def:computation:units}
\begin{enumerate}[label=(\roman*)]
\item The unit of the paper's tables is a curve, one per isogeny class (the optimal one), at a fixed $\ell$. Other units occur: pairs $(E,\ell)$
(Section sec:forms: ``$913{,}244$ working-set pairs $(E,\ell)$''), pairs (curve, bad prime $q$) (Lemma lem:twist and the paragraph after Theorem thm:twist),
triples (curve, $q$, $\ell$). Counts in different units are not comparable.
\item The populations are those of \Cref{conv:computation:populations}, and each filter is part of a statement: CM curves are excluded by
Conjecture conj:euler and Conjecture conj:eis; a reducible $E[\ell]$ moves a curve to the Eisenstein population, where Conjecture conj:eis makes a
different statement; and in the hypothesis population, where Theorem thm:main applies, the data check a theorem rather than test a conjecture.
\item Only one curve per class is counted because $v_\ell(\degphi)$ is attached to the optimal curve; counting every curve of a class would count
one newform several times. For $E[\ell]$ irreducible, $v_\ell$ of the degree and the terms $e_\ell$, $t_\ell$, $h_\ell$ are invariants of the class
(\Cref{ch:moddeg,ch:ellcurves}); the weight $w_3$ is only observed to be one.
\end{enumerate}
\end{definition}

\emph{Two data corrections on the populations.} (a) Section sec:forms: ``The hypothesis of Theorem thm:main fails on the working set for the $1{,}170$
curves with image $\Ns$ at $\ell=3$ \dots; by Proposition prop:Sf there are no other failures among the curves with $E[\ell]$ irreducible''. This holds
for the population without CM. A CM curve with $j=-32768$ or $8000$ and $3\nmid N$ has image $N_s(\F_3)$ and fails the hypothesis at $\ell=3$ too, and
\texttt{galrep} gives it no label: below conductor $10{,}000$ there are $20$ such curves in the tables, $10$ of them optimal. The sentence needs
``without CM'', as Table tab:counts already has. (b) Section sec:boundary, after Proposition prop:ns-structure: ``Among the $90$ optimal curves with
image $\Ns$ of conductor below $10{,}000$ ($28$ of them with $3\mid N$), $74$ have such a prime $p\ge5$ \dots\ and $16$ have it only at $p=2$''. The $90$,
$74$ and $16$ are right for the optimal curves without CM (\Cref{ex:computation:repro-ns}), but $44$ of the $90$ have $3\mid N$ ($14$ with $3\,\|\,N$, $6$
with $9\,\|\,N$, $24$ with $27\mid N$), not $28$; and the proposition allows CM, while the count leaves out the $13$ optimal CM curves with image
$N_s(\F_3)$ of \Cref{ex:computation:mod3-image}.

\begin{definition}[rates]\label{def:computation:rates}
For a population and a bound $B$: the number of curves with $B>0$; the number of violations; the \emph{equality rate}, the number of equalities
divided by the number of curves with $B>0$; the \emph{mean excess}, the average of $v-B$ over the curves with $B>0$; and the \emph{base rate} of an
attribute, its proportion in the population (\Cref{ch:prereq}).
\end{definition}

Three remarks on these. Curves with $B=0$ cannot violate the bound, and including them would inflate every rate; this is why all rates are
taken among positive bounds. An equality rate is a property of a population and a range, not of ``elliptic curves'': at $\ell=3$ it is $75.1\%$
on $N<5000$ and $68.0\%$ on $5000\le N<10000$ (\Cref{ex:computation:repro-uniform}), against the paper's $57.5\%$ on $N\le300000$; degrees grow with
$N$, and with them the room for extra divisibility. The mean excess over $B^{\mathrm{rule}}_3$ is $0.31$, $0.42$ and $0.37$ on the two small ranges and their
union, against the paper's $0.63$ on its working set. Under the hypothesis of Theorem thm:main the excess is not noise: (eq:identity) is an
equality there (\Cref{ch:paper}), the Tamagawa factor of $\ad T$ is trivial at the potentially good additive primes (Lemma lem:tamadd), and at the
potentially multiplicative primes it exceeds $t_\ell$ by the explicit term of Lemma lem:tamfactor; so the excess is the length of the adjoint
Selmer group plus that explicit term (Corollary cor:explicit). The data measure the excess; they do not explain it (Open question 3).

\begin{example}[base rates by Kodaira type]\label{ex:computation:base-rates}
In the base population at $\ell=3$ below conductor $10{,}000$ ($33{,}851$ curves), $3\mid\degphi$ for $76.1\%$. Conditionally on having at least one prime $q\ge5$
of a given type (a curve with several types is counted in several rows):
\begin{center}\small
\begin{tabular}{lrr}
\toprule
type at some $q\ge5$ & curves & share with $3\mid\degphi$\\
\midrule
$\kod{II}$ / $\kod{II}^*$ & $946$ / $418$ & $100\%$ / $100\%$\\
$\kod{IV}$ / $\kod{IV}^*$ & $468$ / $896$ & $100\%$ / $100\%$\\
$\kod{I}_n^*$ & $2{,}233$ & $100\%$\\
$\kod{I}_n$ with $3\mid n$ & $5{,}134$ & $100\%$\\
$\kod{III}$ / $\kod{III}^*$ & $813$ / $807$ & $75.4\%$ / $75.5\%$\\
$\kod{I}_0^*$ & $692$ & $73.1\%$\\
$\kod{I}_n$ with $3\nmid n$ & $31{,}185$ & $74.6\%$\\
\bottomrule
\end{tabular}
\end{center}
The rows at $100\%$ are exactly the types whose term is always positive at $\ell=3$: $e_3\ge1$ for $\kod{II}$ to $\kod{IV}^*$ ($v_3(q-\chi_{-3}(q))\ge1$) and for $\kod{I}_n^*$
($v_3(q^2-1)\ge1$) at $q\ge5$, and $t_3\ge1$ for $\kod{I}_n$ with $3\mid n$. The rows near the base rate are the types whose term is positive only for some $q$:
$\kod{III}$, $\kod{III}^*$ for $q\equiv\pm1\pmod{12}$, $\kod{I}_0^*$ when $q\equiv1\pmod3$ or $3\mid(q+1)^2-a'^2$. A base-rate table is the first thing to compute and
the weakest evidence: it is the shadow of the inequality on a coarse partition. Section sec:iso of the paper refines it cell by cell
(\Cref{ex:computation:repro-iso}).
\end{example}

\begin{remark}[pooling]\label{rem:computation:simpson}
Rates pooled over ranges or over types can mislead. The pooled equality rate at $\ell=3$ below $10{,}000$, $71.0\%$, is the average of the rates of the
two small ranges, $75.1\%$ and $68.0\%$, weighted by their numbers of positive bounds, $8{,}510$ and $11{,}807$. A comparison of the paper's $57.5\%$ on
$N\le300000$ with the book's $71.0\%$ says nothing about a difference in the arithmetic of the curves; it says that the ranges differ, and the
trend is visible inside the book's range already. When the composition of the population changes with the range (more bad primes, larger
degrees), a pooled rate can even move in the opposite direction to every one of its parts (Simpson's paradox). The fair comparison is like with
like: the same statement, the same population, the same range. The paper's equality rate on its out-of-table curves, $52.5\%$ of the $828$
curves with a positive bound at $\ell=3$ (Section sec:oot), is a rate on a population built by a different procedure (\Cref{rem:computation:protocol}).
\end{remark}

\section{The protocol}\label{sec:computation:protocol}

The paper describes a protocol: statements formulated on a working set, written down, tested once on a hold-out, then on further ranges. This
section describes each range, classifies each data statement of the paper as blind-tested, confirmed only, or a theorem checked, explains what
a blind test of a universal inequality establishes, discusses multiple comparisons, and runs the protocol in miniature on the book's data.
In the paper it is the protocol paragraph of Section sec:data, with Sections sec:wild and sec:oot.

\begin{remark}[the four ranges]\label{rem:computation:protocol}
The paper's ranges (\Cref{conv:computation:populations}) are: the working set, $N\le300000$, where every class has determined optimality; the
hold-out, $300000<N\le400000$, the same; the extension, $400000<N\le500000$, where optimality is determined for $359{,}009$ classes only (whether the
extension counts are restricted to those classes is not stated, \Cref{rem:computation:optimality-data}); and the out-of-table curves: $893$ curves (the abstract)
with $500000<N\le4\cdot10^6$; ``we generated curves with prescribed local types'' (``families with a type $\mathrm{IV}$ or $\mathrm{IV}^*$ prime, with two
or three type $\mathrm{II}$ primes, with mixed types, and random curves''), kept when the isogeny class is trivial ``so that the curve is optimal and
$E[\ell]$ is irreducible for every $\ell$'', not CM, and with modular degrees computed by \texttt{ellmoddegree} (Section sec:oot). The vocabulary
(working set, hold-out, pre-registration, blind test, confirmation) is that of \Cref{ch:prereq}.
\end{remark}

\begin{remark}[the status of each data statement]\label{rem:computation:test-status}
We say that a statement is \emph{blind-tested} on a range if it was written down before any quantity of that range entering the statement was
looked at. Generating data before a statement is formulated is necessary for this and not sufficient: the data must also not have been
examined. Reading the protocol paragraph of Section sec:data, with Sections sec:wild, sec:tam and sec:oot, the statements of the paper have the
following status.
\begin{center}\footnotesize
\begin{tabular}{p{2.6cm}p{2.1cm}p{2.3cm}p{1.7cm}p{4.6cm}}
\toprule
statement & formulated on & blind test & confirmed only & remark\\
\midrule
(eq:main) & working set & hold-out (once), then extension & & ``nothing was adjusted after the hold-out test''\\
Conjecture conj:eis & working set & hold-out (once), then extension & & $D\le2$ is a theorem (Proposition prop:no9)\\
(eq:twist) & working set, replacing a weaker form & hold-out and extension, ``tested once'' & & blind only if fixed before the hold-out counts of (eq:main) were seen: the order is not stated; the paper does not count the out-of-table curves as a blind test of it\\
Conjecture conj:euler, Euler terms & working set (minima by type) & out-of-table & hold-out, extension & the terms at $2$ and $3$ ``were fixed afterwards by PARI's Euler factors''\\
weights $w_3$ & working set (minima at $q=3$) & out-of-table, only if it contains curves with $3\mid N$ and potentially good reduction at $3$ & hold-out, extension & ``empirical'' (Section sec:wild)\\
correction $h_3$ & from the exceptions & none & all ranges & ``fitted on, and only confirmed by, data already seen''; $h_3=1$ on $\Ns$ curves is a theorem, the repaired inequality is not\\
III-free form & after the $\Ns$ exceptions & none & all ranges & chosen so as to have no exceptions\\
Tamagawa refinement & from Lemma lem:tamfactor & working set, as a prediction & & tests the Tamagawa term of (eq:identity)\\
Lemmas lem:euler, lem:twist; Prop.\ prop:forms; $h_3=1$ & & & & theorems: the data check computations\\
\bottomrule
\end{tabular}
\end{center}
Two of these entries need a qualification that the text does not settle. The out-of-table curves, ``generated before the rule was formulated'',
are a blind test of the Euler terms only if their modular degrees were not examined while the terms at $2$ and $3$ and the weights $w_3$ were
being fixed; the paper does not say. And for (eq:twist), ``then tested once on the hold-out and the extension'' is a blind test only if
(eq:twist) was written down before the hold-out counts of (eq:main) were seen.
\end{remark}

The decisive sentences are these (Section sec:data): ``Since the hold-out and the extension had already been used for (eq:main) and (eq:twist),
they are a confirmation and not a blind test for Conjecture conj:euler; the out-of-table curves, generated before the rule was formulated, are.
They contain no curve with image $\Ns$, so the correction term $h_3$ was fitted on, and only confirmed by, data already seen; and we do not count
the out-of-table curves as a blind test of (eq:twist).''

\begin{remark}[what a blind test of an inequality establishes]\label{rem:computation:what-is-tested}
A single violation refutes a universal inequality. Zero violations on $n$ fresh curves is evidence whose weight depends on two things: how many
violations would have been expected if the statement were unrelated to the degree (the calibration of \Cref{sec:computation:calibration}), and
whether the fresh curves exercise the statement. The out-of-table families were built around the types $\kod{IV}$, $\kod{IV}^*$, two or three $\kod{II}$ primes,
mixed types and random curves; they contain no $\Ns$ curve, so they cannot test the correction term, and they test the $\kod{III}$ terms only if such
primes occur among them. ``Nothing was adjusted after the hold-out test'' means that no term, weight or population filter of (eq:main) and
Conjecture conj:eis was changed after their hold-out counts were seen; it does not cover the statements formulated later.
\end{remark}

\begin{remark}[multiple comparisons]\label{rem:computation:multiple-comparisons}
\begin{enumerate}[label=(\roman*)]
\item The paper's statements came out of a search: the AI statement describes ``the pre-registered search, which produced the candidate
relation'' and then ``the conjectures and their successive refinements'': from the Tamagawa number to $\Phig_q$, to the twist-symmetrised
weights, to the Euler factors, to the weights $w_3$, to the correction term.
\item If $m$ false candidate statements each survive a working-set check with probability $\pi$, about $m\pi$ survive; a hold-out check with
survival probability $\pi'$ for a false statement leaves about $m\pi\pi'$. This is the whole case for holding data back
(\Cref{exer:computation:multiple}).
\item For a universal inequality, $\pi'$ is tiny when the calibration says that thousands of violations would be expected if the bound were
unrelated to the degree (\Cref{sec:computation:calibration}). The risk of multiple comparisons therefore sits mostly in the \emph{choice of form}
(which weights, which correction, which population filter), where the working set was used repeatedly, and not in the bare existence of a
relation.
\item Features chosen by looking at the data are validated only by data not used to choose them. The weights $w_3$ were defined as observed
minima, so the working set cannot violate them; the III-free form was chosen to have no exceptions. This is the statistician's
multiple-comparisons problem \cite[chapter on multiple testing]{Wasserman04} and the physicist's look-elsewhere effect.
\end{enumerate}
\end{remark}

\begin{example}[the protocol in miniature]\label{ex:computation:toy-protocol}
Freeze the small working set $N<5000$. Following the paper's recipe for $w_3$ (Section sec:wild), take the curves of the base population at $\ell=3$
with $3\mid N$, potentially good reduction at $3$, and no other contribution ($e_3(E,q)=0$ at every bad $q\ne3$, $\kod{III}$ terms included, and $t_3(E,q)=0$
at every bad $q$), and read off the minimum of $v_3(\degphi)$ by the type at $3$:
\begin{center}\small
\begin{tabular}{lccccccc}
\toprule
type at $3$ & $\kod{II}$ & $\kod{II}^*$ & $\kod{III}$ & $\kod{III}^*$ & $\kod{IV}$ & $\kod{IV}^*$ & $\kod{I}_0^*$\\
\midrule
$N<5000$: minimum (curves) & $1$ ($239$) & $2$ ($156$) & $0$ ($221$) & $1$ ($221$) & $1$ ($156$) & $2$ ($239$) & $0$ ($299$)\\
$5000\le N<10000$ & $1$ ($255$) & $2$ ($186$) & $0$ ($251$) & $1$ ($251$) & $1$ ($186$) & $2$ ($255$) & $0$ ($347$)\\
\bottomrule
\end{tabular}
\end{center}
The minima on the small working set are exactly the paper's weights. Test them once on the small hold-out: the minima are the same, and the
rule with these frozen weights has no violation among the $11{,}699$ curves of the base population with $3\mid N$ and $5000\le N<10000$. In all,
$3{,}239$ curves of the base population at $\ell=3$ below $10{,}000$ have a positive weight $w_3(E)$. This replicates a recipe; it discovers nothing,
and the small working set was not blind to the paper: the weights were known in advance. It shows the procedure, nothing more.
\end{example}

\begin{paperlink}\label{pl:computation:protocol}
The protocol paragraph of Section sec:data and the AI statement (``a frozen snapshot of the tables, a working set, a range held back for
one-time tests, and a list of tests fixed before any was run''). \Cref{rem:computation:test-status} classifies the statements; its two
qualifications are questions to the author, not errors.
\end{paperlink}
\section{The calibration argument}\label{sec:computation:calibration}

Section sec:sharp of the paper answers an obvious objection: if $3$ divides most modular degrees anyway, zero violations of a lower bound at $\ell=3$
might mean nothing. Its answer is a calibration, an expected number of violations under a null hypothesis, compared with the observed zero. This
section makes the argument exact, reproduces it on the book's range with two other choices of reference population, says what it shows and
what it does not, and adds the paper's two descriptive statistics of sharpness.

\begin{proposition}[the null model]\label{prop:computation:null-model}
Let a finite population of curves carry a bound $b(E)\in\{0,1,2,\dots\}$ and an observed valuation $v(E)$. Let $H_0$ be the hypothesis that the $v(E)$ are
random variables, each with the same distribution $F$ and independent of the bounds. Let $X$ be the number of violations, $X=\#\{E:b(E)>0,\ v(E)<b(E)\}$, and
$p_E=F(v<b(E))$.
\begin{enumerate}[label=(\roman*)]
\item Under $H_0$, $E[X]=\sum_{b(E)>0}p_E$, with no independence between curves needed. If $\hat F$ is the empirical distribution of $v$ on the curves with
$b=0$, then $\hat E=\sum_{b(E)>0}\hat F(v<b(E))$ is the plug-in estimate of $E[X]$.
\item If moreover the $v(E)$ are independent, then $\operatorname{Var}X=\sum p_E(1-p_E)\le E[X]$ and $P(X=0)=\prod_{b(E)>0}(1-p_E)\le e^{-E[X]}$; and if they are only pairwise
independent, $P(X=0)\le\operatorname{Var}X/E[X]^2\le1/E[X]$ whenever $E[X]>0$.
\item Under $H_0$ the distribution of $v$ over the whole population, bound $0$ and positive together, is also $F$. So a large difference between the
empirical distributions of $v$ on all curves and on the bound-$0$ curves is itself evidence against $H_0$. A permutation test makes this
exact: shuffle the observed values of $v$ among the curves and recount the violations; under $H_0$ (with exchangeable $v$) the observed count is a
draw from the distribution of the recounts.
\end{enumerate}
\end{proposition}

\begin{proof}
(i) $X=\sum_{b(E)>0}1_{\{v(E)<b(E)\}}$, and $E[1_{\{v(E)<b(E)\}}]=P(v(E)<b(E))=p_E$, since $b(E)$ is a fixed number and $v(E)$ has distribution $F$; linearity of expectation
(\Cref{ch:prereq}) needs no independence. (ii) For independent indicators the variance of the sum is the sum of the variances $p_E(1-p_E)$, and this is at
most $\sum p_E$. The event $X=0$ is the intersection of the independent events $v(E)\ge b(E)$, of probabilities $1-p_E\le e^{-p_E}$; multiply. Under pairwise
independence the variance formula still holds (the covariances vanish), and $X=0$ implies $\abs{X-E[X]}\ge E[X]$, so Chebyshev's inequality gives
$P(X=0)\le\operatorname{Var}X/E[X]^2\le E[X]/E[X]^2$. (iii) Every $v(E)$ has distribution $F$ whatever $b(E)$ is. If the values are exchangeable, the vector of
values is equal in distribution to any permutation of it, and the recount after a random permutation has the same distribution as $X$.
\end{proof}

The proposition is elementary; its content is in the choice of $H_0$ and of the reference population, which the next examples vary.

\begin{example}[the paper's numbers]\label{ex:computation:calibration-paper}
Section sec:sharp: ``Among the $366{,}151$ working-set curves with $E[3]$ irreducible, not CM, and right-hand side $0$ in (eq:rule), the prime $3$ does not
divide $\degphi$ for $53.2\%$. If the valuation of the degree were independent of the local data, the $867{,}578$ working-set curves with a positive bound at
$\ell=3$ would produce roughly $605{,}000$ violations (summing over them the proportion of bound-$0$ curves whose valuation is smaller than the bound); there
are none.'' The arithmetic: the curves with $v=0$ alone contribute at least $0.532\cdot867{,}578\approx461{,}550$ to $\hat E$, since $\hat F(v<b)\ge\hat F(v=0)$ for $b\ge1$; the
stated $605{,}000$ means an average $\hat F(v<b)$ of $0.697$ over the bounded curves, close to the book's $0.703$ (\Cref{ex:computation:calibration-small}).
Under independence the standard deviation of $X$ is at most $\sqrt{605000}\approx778$, so the observed $0$ lies about $778$ standard deviations below the
mean, and $P(X=0)\le e^{-605000}$. A number like this is not a meaningful $p$-value: it is far below any doubt about the model itself (independence
across curves is an idealisation). That is why the paper, and \Cref{ch:prereq}, compare $605{,}000$ with $0$ and state no $p$-value. The paper adds:
``This shows that the statement is strong, far from vacuous, not that it is valid; the figure depends on the reference population.''
\end{example}

\begin{example}[the calibration on the book's range]\label{ex:computation:calibration-small}
In the base population at $\ell=3$ below conductor $10{,}000$ ($33{,}851$ curves, bound $B^{\mathrm{rule}}_3$): $13{,}534$ curves have bound $0$, with $v=0,1,2,3,4,5,6$ for
$8{,}084$ ($59.7\%$), $4{,}079$, $1{,}081$, $251$, $33$, $5$, $1$ of them; $20{,}317$ have a positive bound ($1$: $13{,}553$; $2$: $5{,}534$; $3$: $1{,}109$; $4$: $112$; $5$: $9$).
The plug-in expectation is $\hat E=14{,}275$ violations ($70.3\%$ of the bounded curves), with standard deviation $61.5$ under independence; observed:
$0$. By range, with $\hat F$ taken in each range: $\hat E=6{,}091$ on $N<5000$ ($\hat F(v=0)=0.630$) and $8{,}119$ on $5000\le N<10000$ ($\hat F(v=0)=0.566$). At
$\ell=5,7,11,13$ (union of the ranges) $\hat E=5{,}901$, $3{,}476$, $1{,}187$, $698$; observed $0$ each time. Two other reference populations, at $\ell=3$:
\begin{enumerate}[label=(\alph*)]
\item \emph{Stratified.} Split the curves into $16$ strata by conductor (bins of width $2500$) and by the number of prime factors of $N$, capped at $4$, and
take $\hat F$ from the bound-$0$ curves of each stratum: $\hat E=14{,}177$.
\item \emph{Permutation.} Shuffle $v$ among all the curves of each stratum, bound $0$ and positive together, and recount: over $200$ shuffles (Python's
\texttt{random.seed(0)}) the number of violations has mean $7{,}638$, sample standard deviation $49.3$, minimum $7{,}517$ and maximum $7{,}779$.
\end{enumerate}
The reference population changes the expected count by a factor of two and the conclusion not at all. The factor of two has a visible cause:
$v=0$ for $59.7\%$ of the bound-$0$ curves but only for $23.9\%$ of all curves, a difference that \Cref{prop:computation:null-model}(iii) says should not exist
under $H_0$. The bound-$0$ curves are those with no local reason for divisibility, and their degrees are less divisible; the pooled distribution
already contains the curves the bound singles out.
\end{example}

\begin{remark}[what the calibration shows and what it does not]\label{rem:computation:calibration-reading}
It shows that the zero count is not an artefact of $3$ dividing most modular degrees ($76.1\%$ below $10{,}000$, \Cref{ex:computation:base-rates}): the
bound picks out, among all curves, those whose degree is divisible, and a bound unrelated to the degree would fail thousands of times in the book's
range and hundreds of thousands of times in the paper's. It does not show:
\begin{enumerate}[label=(\alph*)]
\item that the inequality holds beyond the data; that is the business of Theorem thm:main under its hypotheses, and of the conjecture outside them;
\item anything about validity on the working set: the rule was fitted there (weights read off minima, terms at $2$ and $3$ fixed after the minima
were seen), so on the working set the calibration measures the \emph{strength} of the statement, as the paper says; its \emph{validity} rests on the
blind tests of \Cref{rem:computation:test-status};
\item a precise level of significance: independence across curves is an idealisation (twists and congruent newforms tie curves together), so the
standard deviations are indicative only;
\item that the null is the right one: ``independent of the local data'' is crude, and conditioning on the conductor and the number of bad primes,
as in \Cref{ex:computation:calibration-small}(a), is a sharper null.
\end{enumerate}
The equality rates and the mean excess are descriptive statistics of the population, not tests (\Cref{def:computation:rates}).
\end{remark}

\begin{example}[sharpness against chance]\label{ex:computation:sharpness}
Under the null of \Cref{prop:computation:null-model} the expected number of equalities among the $20{,}317$ bounded curves at $\ell=3$ below $10{,}000$ is
$\sum\hat F(v=b(E))\approx4{,}548$; observed: $14{,}420$ ($71.0\%$). Over the union of the small ranges, the equality rates of (eq:rule) at $\ell=3,5,7,11,13$ are
$71.0\%$, $84.5\%$, $90.7\%$, $96.5\%$, $96.0\%$, and the mean excesses $0.37$, $0.17$, $0.10$, $0.04$, $0.04$. The paper (working set): ``$v_\ell(\degphi)$ equals the
right-hand side of (eq:rule) in $57.5\%$ of the cases at $\ell=3$, $75.6\%$ at $\ell=5$, $84.1\%$ at $\ell=7$, $91.6\%$ at $\ell=11$ and $92.4\%$ at $\ell=13$, and the
mean excess at $\ell=3$ is $0.63$.'' The ordering by $\ell$ is the same; the levels differ with the range (\Cref{rem:computation:simpson}).
\end{example}

\begin{example}[the odd part of the degree]\label{ex:computation:odd-part}
The paper's last statistic concerns the optimal curves without CM and with no isogeny of odd degree (so that $E[\ell]$ is irreducible for every odd
$\ell$; decided here from the row of the optimal curve in the \texttt{allisog} matrix): how often the odd part of $\degphi$ equals $\prod_\ell\ell^{B^{\mathrm{tw}}_\ell(E)}$, the
product over odd $\ell$. Below conductor $10{,}000$ there are $33{,}137$ such curves. The quotient (odd part of $\degphi$)$/\prod_\ell\ell^{B^{\mathrm{tw}}_\ell}$ is always an
integer, which is (eq:twist) at every odd $\ell$ at once; it is $1$ for $5{,}675$ curves ($17.1\%$; $25.5\%$ of $14{,}666$ curves on $N<5000$ and $10.4\%$ of
$18{,}471$ on $5000\le N<10000$). The most common quotients are $3$ ($5{,}926$ curves), $1$, $5$ ($2{,}978$), $9$ ($2{,}532$), $15$ ($2{,}024$), $7$ ($1{,}918$), $21$
($1{,}088$) and $11$ ($967$). Among the nontrivial quotients, the smallest prime factor is $3$ in $58.7\%$ of the cases and divides $N$ in $47.5\%$. The paper (working
set, $1{,}226{,}216$ curves): equality in $1.85\%$ (``the proportion is range dependent: it is $17.1\%$ below conductor $10{,}000$''), quotients ``most often $9$, $3$,
$15$, $45$, $27$ or $21$'', ``its smallest prime factor is $3$ in $63\%$ of the cases and divides $N$ in $52\%$''. We read the last phrase as ``the smallest prime factor
of the quotient divides $N$''. The local terms account for a sharp part of $\degphi$ and for less and less of the rest as $N$ grows; the rest is the
adjoint Selmer group together with the size of the Petersson norm (\Cref{ch:lfunctions}).
\end{example}
\section{Reproducing the paper's tables below conductor \texorpdfstring{$10{,}000$}{10,000}}\label{sec:computation:repro}

This section reproduces, on all $38{,}042$ optimal curves of conductor below $10{,}000$, every table and every count of the paper's Section sec:data
that the data of this range allow. The paper's table files are not part of the paper's text, so the paper's numbers quoted here are only those
that its text states; they concern its large ranges and are marked as the paper's. The book's counts are reproductions on data that were not
blind (\Cref{sec:computation:intro}). All counts are over the base population at $\ell$ (\Cref{conv:computation:populations}) unless another is named.

\begin{example}[Table tab:uniform in small]\label{ex:computation:repro-uniform}
For each range and $\ell$: the base population; the curves with $B^{\mathrm{rule}}_\ell>0$, the violations of (eq:rule) and the equality rate; the violations of
$B^{\mathrm{nc}}_\ell$ (the rule without the correction); and the curves with $B^{\mathrm{III}}_\ell>0$, the violations and the equality rate of the III-free form.
\begin{center}\footnotesize
\begin{tabular}{llrrrrrrrr}
\toprule
range & $\ell$ & curves & rule $>0$ & viol. & eq. & viol.\ without $h$ & III-free $>0$ & viol. & eq.\\
\midrule
$N<5000$ & 3 & 15,102 & 8,510 & 0 & 75.1\% & 16 & 8,430 & 0 & 75.0\%\\
 & 5 & 16,863 & 3,078 & 0 & 88.4\% & 0 & 3,068 & 0 & 88.3\%\\
 & 7 & 17,107 & 1,644 & 0 & 91.5\% & 0 & 1,642 & 0 & 91.5\%\\
 & 11 & 17,190 & 515 & 0 & 97.5\% & 0 & 515 & 0 & 97.5\%\\
 & 13 & 17,180 & 296 & 0 & 97.3\% & 0 & 296 & 0 & 97.3\%\\
$5000\le N<10000$ & 3 & 18,749 & 11,807 & 0 & 68.0\% & 20 & 11,697 & 0 & 67.6\%\\
 & 5 & 20,456 & 4,398 & 0 & 81.7\% & 0 & 4,368 & 0 & 81.7\%\\
 & 7 & 20,584 & 2,291 & 0 & 90.2\% & 0 & 2,287 & 0 & 90.2\%\\
 & 11 & 20,649 & 741 & 0 & 95.8\% & 0 & 741 & 0 & 95.8\%\\
 & 13 & 20,643 & 432 & 0 & 95.1\% & 0 & 428 & 0 & 95.1\%\\
$N<10000$ & 3 & 33,851 & 20,317 & 0 & 71.0\% & 36 & 20,127 & 0 & 70.7\%\\
 & 5 & 37,319 & 7,476 & 0 & 84.5\% & 0 & 7,436 & 0 & 84.5\%\\
 & 7 & 37,691 & 3,935 & 0 & 90.7\% & 0 & 3,929 & 0 & 90.8\%\\
 & 11 & 37,839 & 1,256 & 0 & 96.5\% & 0 & 1,256 & 0 & 96.5\%\\
 & 13 & 37,823 & 728 & 0 & 96.0\% & 0 & 724 & 0 & 96.0\%\\
\bottomrule
\end{tabular}
\end{center}
The pure GP loop of \Cref{con:computation:table-small} and the Python scripts agree on $N<2000$: \texttt{tablerow(2000, l, 1, dir)} returns
\texttt{[4951, 2517, 0, 2046]}, \texttt{[5622, 915, 0, 845]}, \texttt{[5735, 468, 0, 446]} at $\ell=3,5,7$ (curves, positive bounds, violations,
equalities), and \texttt{tablerow(2000, 3, 0, dir)} returns \texttt{[4951, 2499, 0, 2028]} for the III-free form; the Python counts are the
same. The paper (Section sec:data): ``no violation at any $\ell$'' on the four ranges; the violations without the correction ``are curves with image
$\Ns$''; the III-free form ``holds for every curve of the tables without any correction''; equality $57.5\%$ at $\ell=3$ on the working set. The
book's table agrees in kind on every point: no violation, the uncorrected failures are exactly $\Ns$ curves (\Cref{ex:computation:repro-ns}), and
the III-free form never fails.
\end{example}

\begin{example}[Table tab:counts in small]\label{ex:computation:repro-counts}
For each range and $\ell$: the curves with $B^{\mathrm{main}}_\ell>0$, violations and equality rate; the same for $B^{\mathrm{tw}}_\ell$; and for the Eisenstein
population: its size, the curves with $B^{\mathrm{main}}_\ell>0$, those with a positive deficit $D=B^{\mathrm{main}}_\ell-v$, the violations of the torsion bound
$D\le v_\ell(\#E(\Q)_{\mathrm{tors}})+[\ell=3,\,E[3]\simeq\Z/3\oplus\mu_3]$ and of the class bound ($D\le$ the number of curves of the class with a rational point of order
$\ell$), and the curves with $D=2$.
\begin{center}\scriptsize\setlength{\tabcolsep}{3pt}
\setlength{\tabcolsep}{2.5pt}\begin{tabular}{llrrrrrrrrrrrr}
\toprule
range & $\ell$ & main $>0$ & viol. & eq. & twist $>0$ & viol. & eq. & Eis. & main $>0$ & $D>0$ & viol. & viol.\ (class) & $D=2$\\
\midrule
$N<5000$ & 3 & 5,908 & 0 & 59.0\% & 6,996 & 0 & 61.3\% & 2,096 & 1,725 & 284 & 0 & 0 & 11\\
 & 5 & 2,478 & 0 & 86.5\% & 2,681 & 0 & 87.2\% & 335 & 199 & 45 & 0 & 0 & 0\\
 & 7 & 1,351 & 0 & 89.8\% & 1,452 & 0 & 90.6\% & 91 & 35 & 9 & 0 & 0 & 0\\
 & 11 & 450 & 0 & 96.4\% & 472 & 0 & 97.2\% & 8 & 0 & 0 & 0 & 0 & 0\\
 & 13 & 272 & 0 & 96.3\% & 286 & 0 & 97.2\% & 18 & 2 & 0 & 0 & 0 & 0\\
$5000\le N<10000$ & 3 & 8,069 & 0 & 49.1\% & 9,545 & 0 & 51.6\% & 1,906 & 1,624 & 147 & 0 & 0 & 3\\
 & 5 & 3,307 & 0 & 76.6\% & 3,636 & 0 & 78.6\% & 199 & 139 & 8 & 0 & 0 & 0\\
 & 7 & 1,719 & 0 & 86.9\% & 1,894 & 0 & 87.8\% & 71 & 32 & 3 & 0 & 0 & 0\\
 & 11 & 579 & 0 & 94.3\% & 631 & 0 & 94.9\% & 6 & 0 & 0 & 0 & 0 & 0\\
 & 13 & 369 & 0 & 93.8\% & 389 & 0 & 94.1\% & 12 & 2 & 0 & 0 & 0 & 0\\
$N<10000$ & 3 & 13,977 & 0 & 53.3\% & 16,541 & 0 & 55.7\% & 4,002 & 3,349 & 431 & 0 & 0 & 14\\
 & 5 & 5,785 & 0 & 80.8\% & 6,317 & 0 & 82.2\% & 534 & 338 & 53 & 0 & 0 & 0\\
 & 7 & 3,070 & 0 & 88.1\% & 3,346 & 0 & 89.0\% & 162 & 67 & 12 & 0 & 0 & 0\\
 & 11 & 1,029 & 0 & 95.2\% & 1,103 & 0 & 95.9\% & 14 & 0 & 0 & 0 & 0 & 0\\
 & 13 & 641 & 0 & 94.9\% & 675 & 0 & 95.4\% & 30 & 4 & 0 & 0 & 0 & 0\\
\bottomrule
\end{tabular}
\end{center}
The paper (Section sec:forms, working set): equality in (eq:twist) for ``$39.4\%$ at $\ell=3$, $71.1\%$ at $\ell=5$, $81.2\%$ at $\ell=7$ and $91.5\%$ at $\ell=13$, and in
(eq:main) for $36.9\%$, $69.7\%$, $80.4\%$ and $91.2\%$''; the rule ``bounds $867{,}578$ curves instead of $705{,}563$ at $\ell=3$'' (in the book's range $20{,}317$ instead of
$16{,}541$). The ordering is the same in both: main below twist below rule, for the number of bounded curves and for the equality rates.
\end{example}

\begin{example}[Table tab:ns in small]\label{ex:computation:repro-ns}
The optimal curves without CM and with image $\Ns$, the violations of $B^{\mathrm{nc}}_3$ (all types, no correction), and the residues modulo $12$ of the
$\kod{III}$-type primes $q\ne3$ with $e_3(E,q)\ge1$ of the violating curves (or $q=2$):
\begin{center}\footnotesize\setlength{\tabcolsep}{4pt}
\begin{tabular}{lrp{1.2cm}p{1.4cm}p{1.2cm}rp{2.4cm}p{1cm}}
\toprule
range & $\Ns$ curves & positive $\kod{III}$ term & violate $B^{\mathrm{nc}}_3$ & other curves violating & \texttt{ns3} $=1$ & residues of the $\kod{III}$ primes & short\-fall\\
\midrule
$N<5000$ & 30 & 30 & 16 (53\%) & 0 & 30 & $1$: 4, $11$: 8, $q=2$: 4 & always 1\\
$5000\le N<10000$ & 60 & 60 & 20 (33\%) & 0 & 60 & $1$: 8, $11$: 12 & always 1\\
$N<10000$ & 90 & 90 & 36 (40\%) & 0 & 90 & $1$: 12, $11$: 20, $q=2$: 4 & always 1\\
\bottomrule
\end{tabular}
\end{center}
Every one of the $36$ violators has a prime $p\equiv-1\pmod3$ with $t_3(E,p)\ge1$, and $12$ of them have $3\mid N$. Of the $90$ curves, $74$ have a prime $q\ge5$ of type
$\kod{III}$ or $\kod{III}^*$ with $e_3\ge1$ ($22$ and $52$ on the two ranges) and the other $16$ get their positive $\kod{III}$ term only at $q=2$, from PARI's factor; $44$
have $3\mid N$. Proposition prop:ns-structure (proved in \Cref{ch:galrep}) explains the first two observations; the data check them.

\emph{Which inequality.} Section sec:anomaly says that Table tab:ns ``records what happens to the inequality (eq:thm), with all types included and no
correction''. The book counts violations of $B^{\mathrm{nc}}_3$, the rule with every term and no correction. Read literally, (eq:thm) has terms at $q\ne\ell$ only; with
that reading there are $24$ violators below $10{,}000$ ($8$ and $16$ on the two ranges) instead of $36$, and still no curve without image $\Ns$ violates.
The revised paper records the two counts (``below conductor $10{,}000$ the two readings give $36$ and $24$ violations'') and attributes the difference
to ``whether the weights at $q=3$ are included''. In fact, below $10{,}000$ all twelve curves counted by one reading and not the other (\curve{726a1},
\curve{726f1}, \curve{2178b1}, \curve{2178h1}, \curve{3840h1}, \curve{3840k1}, \curve{3840o1}, \curve{3840r1}, \curve{5070h1}, \curve{5070p1}, \curve{5808y1}, \curve{5808z1})
differ by the Tamagawa term $t_3(E,3)=1$ of a type $\kod{I}_n$ or $\kod{I}_n^*$ at $3$, and none by a weight $w_3$. The paper (D): $364$, $68$, $43$ and $0$ violations on its four ranges,
$31\%$, $23\%$, $18\%$ of the $\Ns$ curves. The falling rate with $N$ is visible inside the book's range already ($53\%$, then $33\%$).
\end{example}

\begin{example}[Section sec:iso in small]\label{ex:computation:repro-iso}
Curves of the base population at $\ell$ with exactly one additive prime $q$, with $q\ge5$, $q\ne\ell$, and no multiplicative prime with $\ell\mid n$: the minimum of
$v_\ell(\degphi)$ by the type at $q$ and the value of $e=e_\ell(E,q)$, with the number of curves. At $\ell=3$:
\begin{center}\footnotesize
\begin{tabular}{lll}
\toprule
type & $e$: minimum (curves) & \\
\midrule
$\kod{II}$ & $1$: $1$ ($264$); $2$: $2$ ($27$) & $\kod{II}^*$\quad $1$: $1$ ($135$); $2$: $2$ ($11$)\\
$\kod{IV}$ & $1$: $1$ ($150$); $2$: $2$ ($11$) & $\kod{IV}^*$\quad $1$: $1$ ($249$); $2$: $2$ ($27$)\\
$\kod{III}$ & $0$: $0$ ($193$); $1$: $1$ ($44$); $2$: $2$ ($2$) & $\kod{III}^*$\quad $0$: $0$ ($191$); $1$: $1$ ($44$); $2$: $2$ ($2$)\\
$\kod{I}_0^*$ & $0$: $0$ ($146$); $1$: $1$ ($42$); $2$: $2$ ($88$); $3$: $3$ ($50$) & $\kod{I}_n^*$\quad $1$: $1$ ($813$); $2$: $2$ ($57$); $3$: $3$ ($3$)\\
\bottomrule
\end{tabular}
\end{center}
In every $\kod{I}_n^*$ cell the excess $v_3(\degphi)-e-t_3$ has minimum $0$, also cell by cell in $(e,t_3)$; that is, the excess over $e$ has minimum exactly
$t_3(E,q)=v_3(n)$, as the paper says. At $\ell=5$ every cell's minimum equals $e$: $\kod{I}_0^*$ $e=0,1,2$ ($94$, $75$, $10$ curves), $\kod{II}$ $0,1$ ($210$, $12$), $\kod{II}^*$ ($94$,
$5$), $\kod{III}$ ($152$, $11$), $\kod{III}^*$ ($150$, $11$), $\kod{IV}$ ($95$, $5$), $\kod{IV}^*$ ($209$, $12$), $\kod{I}_n^*$ ($427$, $141$); and at $\ell=7$ likewise, with $e\le1$ (the largest
cell $\kod{I}_n^*$ with $e=0$, $806$ curves). \emph{Multiplicative primes alone:} for the curves with no additive prime and exactly one multiplicative prime $q$ with
$\ell\mid n$, $\min(v_\ell(\degphi)-v_\ell(n))=0$ for every value of $v_\ell(q^2-1)$: at $\ell=3$, $4{,}367$ curves, with $v_3(q^2-1)=1,2,3,4$ for $3{,}094$, $268$, $36$, $3$
curves and $q=3=\ell$ for $966$; at $\ell=5$, $2{,}588$ curves. \emph{Two primes:} the curves of the base population at $\ell=3$ whose only contributions are two odd
primes of type $\kod{II}$ or $\kod{II}^*$ are $8$ (\curve{2450d1}, \curve{3675k1}, \curve{6050bg1}, \curve{7350e1}, \curve{7350ba1}, \curve{7350bw1}, \curve{8281h1}, \curve{9800bj1}), and
$9\mid\degphi$ for all $8$. The paper (D, working set): the same pattern in every cell; the cells with $e=0$ have $5{,}534$ ($\kod{III}$) and $5{,}394$ ($\kod{I}_0^*$)
curves; $\kod{I}_0^*$ runs up to $e=6$; $2{,}502$ curves with two primes of type $\kod{II}$ or $\kod{II}^*$. One small point of wording: the paper's ``for every value of
$v_\ell(q^2-1)$ up to $6$'' includes, at $\ell=3$, the value $0$ only through $q=\ell=3$, where $v_3(q^2-1)$ is not the relevant quantity.
\end{example}

\begin{example}[Section sec:tam in small]\label{ex:computation:repro-tam}
The hypothesis population at $\ell$ (\Cref{conv:computation:populations}), restricted to the curves with at least one prime $q$ of type $\kod{I}_n$ with $\ell\mid n$ and
$q\equiv1\pmod\ell$; the extra term $\sum_q\min(v_\ell(n),v_\ell(q-1),v_\ell(c(u)))$ over such primes, from \texttt{ellinit(E, O(q\^{}30)).tate} as in
\Cref{ex:computation:tate-parameter}; and the refined bound $B^{\mathrm{rule}}_\ell+$ (extra term).
\begin{center}\footnotesize
\begin{tabular}{lrrrrr}
\toprule
& curves & positive extra term & violations & equality before & equality after\\
\midrule
$\ell=3$, type $\kod{I}_n$ & 989 & 179 & 0 & 550 (55.6\%) & 662 (66.9\%)\\
$\ell=3$, with $\kod{I}_n^*$ & 1,050 & 186 & 0 & 596 (56.8\%) & 714 (68.0\%)\\
$\ell=5$, type $\kod{I}_n$ & 246 & 24 & 0 & 197 (80.1\%) & 214 (87.0\%)\\
$\ell=5$, with $\kod{I}_n^*$ & 252 & 24 & 0 & 203 (80.6\%) & 220 (87.3\%)\\
$\ell=7$ & 18 & 1 & 0 & 17 & 18\\
\bottomrule
\end{tabular}
\end{center}
The extra term is never more than $1$ in this range. The rows ``with $\kod{I}_n^*$'' also use the odd primes of type $\kod{I}_n^*$ with $\ell\mid n$ and $q\equiv1\pmod\ell$,
whose Tate parameter is that of the twist by $q^*$, which is multiplicative at $q$ (\Cref{ex:computation:tate-parameter}); the prime $2$ never qualifies. The first
curves with a positive extra term are \curve{791a1} ($q=7$), \curve{805b1} ($7$, $\bar u=1$), \curve{805c1}, \curve{805d1}, \curve{874a1} ($19$, $\bar u=12$), \curve{874b1}
($19$, $\bar u=18$) at $\ell=3$; \curve{4136e1}, \curve{4136f1} ($11$, $\bar u=10$), \curve{5673a1} ($61$), \curve{6507a1}, \curve{6507b1} ($241$) at $\ell=5$; and \curve{4171a1}
($43$, $\bar u=1$) at $\ell=7$. The paper (D, working set, type $\kod{I}_n$ only): at $\ell=3$, $46{,}222$ curves, $12{,}095$ with a positive extra term, equality from
$45.0\%$ to $60.6\%$; at $\ell=5$, $10{,}940$, $1{,}646$, from $67.0\%$ to $78.7\%$; and, including $\kod{I}_n^*$ below conductor $10{,}000$, ``$1{,}050$ such curves at $\ell=3$, with
equality rising from $56.8\%$ to $68.0\%$, and $252$ at $\ell=5$, with equality rising from $203$ to $220$ curves'', which the table reproduces exactly.
\end{example}

\begin{example}[Section sec:eis in small]\label{ex:computation:repro-eis}
At $\ell=3$: $4{,}002$ Eisenstein curves, $3{,}349$ with $B^{\mathrm{main}}_3>0$, $431$ with $D>0$, $D=2$ for $14$. No curve without a rational point of order $3$ has $D>0$, and
the torsion bound is attained with $D>0$ in $427$ cases. $D=2$ occurs for twelve curves with image \texttt{3Cs.1.1} (\curve{14a1}, \curve{26a1}, \curve{35a1},
\curve{38a1}, \curve{77b1}, \curve{158d1}, \curve{278b1}, \curve{326c1}, \curve{485a1}, \curve{4385b1}, \curve{7094c1}, \curve{8027a1}), all of rank $0$, with torsion $\Z/3$
except \curve{14a1} (torsion $\Z/6$), and for two curves with a point of order $9$ (\curve{714i1}, \curve{6942n1}, image \texttt{3B.1.1}). Of the $49$ optimal curves
without CM with image \texttt{3Cs.1.1}, none has a point of order $9$ (torsion $\Z/3$ for $48$ and $\Z/6$ for \curve{14a1}); the four optimal curves with a point of
order $9$ (\curve{714i1}, \curve{1482l1}, \curve{1554n1}, \curve{6942n1}) have image \texttt{3B.1.1}. That the two sources of a deficit $2$ never combine is a theorem
(Proposition prop:no9). At $\ell=5$ and $7$, $D\le1$, and $D>0$ only with a rational point of order $\ell$ ($53$ and $12$ curves). Measured against (eq:twist), the
deficit $B^{\mathrm{tw}}_\ell-v$ exceeds the torsion bound for $180$ curves at $\ell=3$, $35$ at $\ell=5$ and $10$ at $\ell=7$. The paper (D, working set): $75{,}430$
Eisenstein curves, $67{,}845$ with a positive bound, $2{,}777$ with $D>0$, $D=2$ for $22$ (``seventeen curves, all of rank $0$, with torsion $\Z/3$ or $\Z/6$, for
instance 14a1'', and five with a point of order $9$); the deficit against (eq:twist) exceeds the bound for $1{,}148$ curves at $\ell=3$ and $141$ at $\ell=5$. The
book's data agree with the revised wording; the original version's ``all of rank $0$ and torsion $\Z/3$'' was contradicted by \curve{14a1}.
\end{example}

\begin{example}[the twist pairs after Theorem thm:twist]\label{ex:computation:repro-twist}
For each optimal curve without CM below $10{,}000$ and each odd $q$ with $q^2\,\|\,N$ in case (a) (type $\kod{I}_n^*$), (b) (type $\kod{I}_0^*$) or (c) (potentially good, $q\ge5$),
let $E_q$ be the twist by $q^*$, found in \texttt{allcurves} through \texttt{ellminimalmodel}, and compare $v_\ell$ of the two optimal degrees for each
$\ell\in\{3,5,7,11,13\}$ with $\ell\ne q$ and $\ell^2\nmid N$. Pairs (curve, $q$): (a) $6{,}348$, (b) $2{,}183$, (c) $5{,}150$; every twist was found in the tables.
Triples (curve, $q$, $\ell$): (a) $24{,}373$, (b) $8{,}348$, (c) $19{,}554$. In all of them $v_\ell(\degphi)\ge v_\ell(\deg\varphi_{E_q})$; the inequality is strict in $2{,}779$
triples of case (a), in $1{,}308$ of case (b) and in none of case (c). Examples in case (a) at $\ell=3$: \curve{121d1} against \curve{11a1} ($\degphi$ $24$ against $1$),
\curve{75b1} against \curve{15a1} ($6$ against $1$), \curve{150c1} against \curve{30a1} ($48$ against $2$), \curve{539c1} against \curve{77c1} ($288$ against $6$),
\curve{980i1} against \curve{140b1} ($2880$ against $60$); in case (b), \curve{275b1} against \curve{11a1} at $\ell=7$ ($28$ against $1$) and \curve{539d1} against \curve{11a1} at
$\ell=3$ ($72$ against $1$). When moreover $\ell\nmid N$ and $E[\ell]$ is irreducible, the difference equals the Euler term $e_\ell(E,q)$ in all $19{,}691$ triples of case (a)
and all $6{,}701$ of case (b), as Theorem thm:quant predicts (\Cref{ch:congruence}). The paper (D, working set): $157{,}938$ and $292{,}580$ pairs, ``with no exception,
with equality in every pair of case (c)''; its revised text also notes the strict inequalities in case (a). The original version's ``strict inequality
only in case (b)'' is contradicted already below $10{,}000$.
\end{example}

\begin{example}[a twist relation at $q=\ell$, observed]\label{ex:computation:twist-at-ell}
The weights $w_3$ of Conjecture conj:euler have a pattern: under the twist by $\Q(\sqrt{-3})$, which exchanges $\kod{II}\leftrightarrow\kod{IV}^*$,
$\kod{III}\leftrightarrow\kod{III}^*$, $\kod{IV}\leftrightarrow\kod{II}^*$ at $3$ (\Cref{ex:computation:twists}), the starred type always has weight one more than its partner ($2=1+1$,
$1=0+1$, $2=1+1$). The book's data show the same phenomenon at every $\ell$, as a relation between the degrees of a curve and of its twist at
$q=\ell$. For each optimal curve $E$ without CM below conductor $10{,}000$ with $\ell\mid N$, $\ell\in\{3,5,7,11,13\}$, let $E^{(\ell)}$ be its twist by $\chi_{\ell^*}$ and
compare $v_\ell(\degphi)$ with $v_\ell$ of the degree of the optimal curve of the class of $E^{(\ell)}$:
\begin{enumerate}[label=(\alph*)]
\item if $E$ has type $\kod{II}^*$, $\kod{III}^*$ or $\kod{IV}^*$ at $\ell$, the difference is exactly $1$ in all $4{,}412$ cases ($2{,}141$, $1{,}210$, $659$, $239$, $163$ at
$\ell=3,5,7,11,13$; $4{,}251$ of them with $E[\ell]$ irreducible);
\item if $E$ has type $\kod{I}_n^*$ at $\ell$, the difference is $0$ in all cases at $\ell=3,5,7$;
\item if $E$ has type $\kod{I}_0^*$ at $\ell$ and $E[\ell]$ is irreducible, the difference is $v_\ell\bigl((\ell+1)^2-a_\ell(E^{(\ell)})^2\bigr)$ in all $1{,}915$ cases (it fails in $13$ of
the cases with $E[\ell]$ reducible, for instance \curve{126a1} at $\ell=3$ and \curve{275b1} at $\ell=5$).
\end{enumerate}
In (b) and (c) the difference is what the formula for $e_\ell(E,q)$ of Lemma lem:euler gives when $q$ is replaced by $\ell$ ($v_\ell(\ell^2-1)=0$, and
$v_\ell(\ell-1)=0$); in (a) it is not, since at $q\ne\ell$ a potentially good pair has difference $0$ (\Cref{ex:computation:repro-twist}). At $\ell\ge5$, (a) suggests a
weight $1$ for the types $\kod{II}^*$, $\kod{III}^*$, $\kod{IV}^*$ at $q=\ell$, which the paper does not use (``at $q=\ell\ge5$ no term is needed''); added to (eq:rule), it gives no
violation below $10{,}000$ at $\ell=5,7,11,13$ and is attained by most of these curves. These are observations made on the whole of the book's range at
once, with no hold-out, and no explanation is offered here; they bear on the paper's first open question, which asks for the weights at
$q=\ell$.
\end{example}

\begin{example}[further checks]\label{ex:computation:repro-further}
Below conductor $10{,}000$, on the optimal curves:
\begin{enumerate}[label=(\alph*)]
\item Lemma lem:euler and Lemma lem:twist: \Cref{ex:computation:euler-by-hand,ex:computation:twists}.
\item Proposition prop:shapes at $q=2,3$: the cubic factors, \Cref{ex:computation:wild-two}.
\item Proposition prop:forms: all $3{,}226$ optimal curves with type $\kod{IV}$ or $\kod{IV}^*$ at $2$ ($3{,}194$ of them without CM) have $P_2(X)=1+2X$. The paper: ``The
previous version excepted the types $\mathrm{IV}$, $\mathrm{IV}^*$ at $2$ and checked $e_3(E,2)\ge1$ on the $92{,}812$ such curves of the working set, where PARI's
factor at $2$ is $1+2X$ for every one of them; Lemma lem:IV shows that it must be.'' The data check a theorem.
\item Section sec:two: for the $25{,}086$ optimal curves without CM and without a label \texttt{2B} or \texttt{2Cs} (no rational $2$-isogeny), the inequality
$v_2(\degphi)\ge\sum_{q\,\|\,N}v_2(n_q)+\#\{q:\ \text{type }\kod{III},\kod{III}^*,\kod{I}_0^*\text{ or }\kod{I}_n^*\}$ holds with no exception, while counting $\kod{I}_0^*$ with weight $2$ fails
$6$ times (\curve{99d1}, \curve{931b1}, \curve{2299d1}, \curve{2475h1}, \curve{5819a1}, \curve{5931b1}), as the paper's $123$, $26$ and $20$ exceptions on its ranges lead one to expect.
\item The optimal CM curves of odd modular degree are \curve{27a1}, \curve{32a1}, \curve{36a1}, \curve{49a1} (degree $1$) and \curve{243b1} (degree $9$); \curve{27a1} and
\curve{243b1} have torsion $\Z/3$ and conductor neither prime nor with rational $2$-torsion, so they fall under the CM case of the theorem of Calegari and Emerton
quoted in Section sec:two (\Cref{ch:moddeg} states it).
\end{enumerate}
\end{example}

\begin{remark}[what agrees and what does not]\label{rem:computation:comparison}
The book's range is about fifty times smaller in conductor than the paper's working set, and every rate that depends on $N$ differs: the equality rates,
the share of $\Ns$ curves that fail without the correction, the odd-part statistic, the mean excess. All of them move with $N$ inside the book's range
already, in the direction of the paper's values. Every qualitative claim of Section sec:data reproduces: no violation of (eq:rule), (eq:main),
(eq:twist) or of Conjecture conj:eis; the uncorrected failures are exactly $\Ns$ curves and fail by exactly one; the minima equal $e_\ell$ cell by cell; the
weights $w_3$ are the minima at $q=3$; the Tamagawa refinement has no exception and its counts below $10{,}000$ are the paper's; the strict twist
inequalities occur in cases (a) and (b). The sentences that do not reproduce as written are the data corrections recorded in
\Cref{sec:computation:populations} and \Cref{ex:computation:wild-two,ex:computation:repro-ns}.
\end{remark}

\section{Reproducibility}\label{sec:computation:reproducibility}

This last section describes the repository that the paper's appendix names, prints the book's own scripts, and states a reproducibility checklist
applied to the book's runs. In the paper it is the appendix and the last sentences of the AI statement.

\begin{remark}[the paper's repository]\label{rem:computation:paper-repository}
The appendix says: ``All scripts, the pinned data commit and the result tables are in the repository
\url{https://github.com/dsr001001/mcptrial} (directory \texttt{ecmine}).'' It names \texttt{scripts/step10\_sym2\_euler.py} for the
symmetric-square Euler factors, \texttt{scripts/step10\_rule.py} and \texttt{scripts/step11\_tw\_boundary.py} for the rule on the whole ranges
(``correction term, image statistics''), \texttt{scripts/step12\_tamagawa.py} for the Tamagawa factors, \texttt{scripts/conj\_check.py} for the two
corollary forms, \texttt{scripts/euler\_rule.gp} for the functions of \Cref{con:computation:rule-listing}, and \texttt{twist\_lemma\_check.gp}; it adds
that ``each runs on the full tables in minutes to hours'' and that ``The pipeline was rerun from a clean clone and reproduced every count of Table
tab:counts.'' This book has not used the repository: the description above is the appendix's, and the scripts below are the book's own.
\end{remark}

\begin{construction}[the script \texttt{sessions.gp}]\label{con:computation:sessions}
The functions used in the sessions of \Cref{sec:computation:invariants,sec:computation:sessions}; \texttt{rule\_terms} repeats the loop of
\texttt{rule} and records each row. \texttt{eulercheck} returned \texttt{[]} on every optimal curve below $10{,}000$ (\Cref{ex:computation:euler-by-hand}).
\begin{lstlisting}[language=GP,basicstyle=\ttfamily\scriptsize]
\\ sessions.gp: per-prime breakdown of the rule, Tate data, and the twist check.
\\ rule_terms(E, l, iii): one row [q, code, P_q(x), P_q(q^-2), e, w, t] per bad prime q.
rule_terms(E, l, {iii = 1}) =
{ my(R = ellglobalred(E), P = R[4][,1], K = R[5], S = lfunsympow(E, 2), M = []);
  for(i = 1, #P,
    my(q = P[i], k = K[i][2], F = 1/lfuneuler(S, q), r = subst(F, x, 1/q^2),
       e = 0, w = 0, t = 0, vj = valuation(E.j, q));
    if(q != l,
      my(rr = if(k >= 5, r / (1 - 1/q^2), r));
      if(abs(k) != 3 || iii, e = valuation(rr, l)),
      if(l == 3, w = if(k == 4 || k == 2, 1, if(k == -4 || k == -2, 2, if(k == -3, 1, 0)))));
    if(vj < 0, t = valuation(-vj, l));
    M = concat(M, [[q, k, F, r, e, w, t]]));
  M; }
\\ tatedata(E, q, l): [n, q_E, ubar, k], k the largest k <= v_l(q-1) with ubar an l^k-th power
tatedata(E, q, l) =
{ my(Eq = ellinit(E, O(q^30)), qE = Eq.tate[3], n = valuation(qE, q),
     u = lift(qE / q^n + O(q)), v = valuation(q - 1, l), k = 0);
  while(k < v && Mod(u, q)^((q - 1) / l^(k + 1)) == 1, k++);
  [n, qE, u, k]; }
\\ twistcheck(E): for each odd bad prime q, [q, code of E at q, code of the twist by q* at q, conductor of the twist]
twistcheck(E) =
{ my(P = ellglobalred(E)[4][,1], V = []);
  for(i = 1, #P, my(q = P[i]); if(q == 2, next);
    my(qs = if(q % 4 == 1, q, -q), Et = ellinit(elltwist(E, qs)));
    V = concat(V, [[q, elllocalred(E, q)[2], elllocalred(Et, q)[2], ellglobalred(Et)[1]]]));
  V; }
\\ euler_by_type(E, q): L_q(A_f,1)^(-1) from the Kodaira type, bad prime q >= 5 (Lemma lem:euler)
euler_by_type(E, q) =
{ my(k = elllocalred(E, q)[2]);
  if(k >= 5 || k <= -5, return(1 - 1/q^2));
  if(abs(k) == 2 || abs(k) == 4, return(1 - kronecker(-3, q)/q));
  if(abs(k) == 3, return(1 - kronecker(-4, q)/q));
  if(k == -1, my(a = ellap(ellinit(elltwist(E, if(q % 4 == 1, q, -q))), q));
    return((q - 1)*((q + 1)^2 - a^2)/q^3));
  error("not a bad prime"); }
\\ eulercheck(E): the bad primes q >= 5 where PARI's factor and the table differ (expected: [])
eulercheck(E) =
{ my(S = lfunsympow(E, 2), P = ellglobalred(E)[4][,1], bad = []);
  for(i = 1, #P, my(q = P[i]); if(q < 5, next);
    if(subst(1/lfuneuler(S, q), x, 1/q^2) != euler_by_type(E, q), bad = concat(bad, q)));
  bad; }
\end{lstlisting}
\end{construction}

\begin{construction}[the script \texttt{table\_small.gp}]\label{con:computation:table-small}
One row of \Cref{ex:computation:repro-uniform} in pure GP, from the files \texttt{degphi} and \texttt{galrep} in the directory \texttt{dir}; it needs
\texttt{euler\_rule.gp} for \texttt{ns3}. \texttt{tablerow(Nmax, l, iii, dir)} returns [curves, positive bounds, violations, equalities] for the base
population at $\ell$ with $N<$ \texttt{Nmax}. Below $2{,}000$ each call takes a few seconds.
\begin{lstlisting}[language=GP,basicstyle=\ttfamily\scriptsize]
\\ table_small.gp: the paper's rule on Cremona's optimal curves of conductor < Nmax.
\\ Needs euler_rule.gp (the paper's appendix) for ns3; reads the ecdata files degphi.* and galrep.*.
rulebound(E, l, iii) =
{ my(R = ellglobalred(E), P = R[4][,1], K = R[5], S = lfunsympow(E, 2), b = 0, j = E.j);
  for(i = 1, #P,
    my(q = P[i], k = K[i][2], e = 0, t = 0, vj = valuation(j, q));
    if(q != l,
      my(r = subst(1/lfuneuler(S, q), x, 1/q^2));
      if(k >= 5, r /= (1 - 1/q^2));
      if(abs(k) != 3 || iii, e = valuation(r, l)),
      if(l == 3, e = if(k == 4 || k == 2, 1, if(k == -4 || k == -2, 2, if(k == -3, 1, 0)))));
    if(vj < 0, t = valuation(-vj, l));
    b += e + t);
  if(iii && l == 3, b -= ns3(E));
  b; }
\\ the thirteen rational CM j-invariants
{ CMJ = Set([0, 1728, -3375, 8000, -32768, 54000, 287496, -884736, -12288000, 16581375,
         -884736000, -147197952000, -262537412640768000]); }
\\ galrep: a Map from a curve label to its list of non-surjective image labels
readgal(file) =
{ my(M = Map()); foreach(readstr(file), s, my(w = strsplit(s, " "));
    mapput(M, w[1], w[2..#w])); M; }
strprefix(s, p) = #s >= #p && Vecsmall(s)[1..#p] == Vecsmall(p);
\\ E[l] reducible iff a label begins with "lB" or "lCs" (l followed by a letter, so 37B is not 3B)
reducible(G, lab, l) =
{ my(L); if(!mapisdefined(G, lab, &L), return(0));
  for(i = 1, #L, my(s = L[i]);
    if(strprefix(s, Str(l, "B")) || strprefix(s, Str(l, "Cs")), return(1)));
  0; }
\\ one row of the table: [curves, positive bounds, violations, equalities]
tablerow(Nmax, l, iii, dir) =
{ my(G = readgal(Str(dir, "/galrep.00000-09999")), c = [0, 0, 0, 0]);
  foreach(readstr(Str(dir, "/degphi.00000-09999")), s,
    my(w = strsplit(s, " "), N = eval(w[1]));
    if(N >= Nmax, next);
    my(E = ellinit(eval(w[6])), deg = eval(w[4]), lab);
    if(setsearch(CMJ, E.j), next);
    \\ the label of the optimal curve is N class number, except in class 990h (990h3)
    lab = Str(w[1], w[2], if(N == 990 && w[2] == "h", 3, w[3]));
    if(reducible(G, lab, l), next);
    my(b = rulebound(E, l, iii), v = valuation(deg, l));
    c[1]++; if(b > 0, c[2]++; if(v < b, c[3]++); if(v == b, c[4]++)));
  c; }
\end{lstlisting}
\end{construction}

\begin{remark}[the book's Python scripts]\label{rem:computation:book-scripts}
The counts of \Cref{sec:computation:populations,sec:computation:calibration,sec:computation:repro} come from five Python scripts using \texttt{cypari2}.
\begin{enumerate}[label=(\roman*)]
\item \texttt{compute.py} writes one JSON record per optimal curve: the label (looked up in \texttt{allcurves} from the equation of \texttt{degphi}), $N$,
$\degphi$, $j$, CM, the \texttt{galrep} labels, \texttt{ns3}, rank, torsion, the row of the isogeny matrix and the torsion orders of the class, and for each bad
prime $q$: $f_q$, the code, $c_q$, $P_q$, $P_q(q^{-2})$, $a_q$, $v_q(j)$, the value of \texttt{euler\_by\_type} for $q\ge5$, and for odd $q$ the code, conductor exponent and
minimal model of the twist by $q^*$. It calls one GP function per curve and runs in about $13$ seconds on all $38{,}042$ curves. Its core:
\begin{lstlisting}[basicstyle=\ttfamily\scriptsize]
pari = cypari2.Pari(); pari.allocatemem(64*10**6, 2*10**9, silent=True)
for f in ['euler_rule.gp', 'sessions.gp', 'compute.gp']: pari(f'read("{f}")')
locd = pari('locdata')          # [N, j, ns3, per-prime data], one GP call per curve
for line in open(D + 'degphi.00000-09999'):
    w = line.split(); a = eval(w[5]); lab = bya[w[5]]   # bya: equation -> label, from allcurves
    N, j, ns3, V = locd(pari.ellinit(a))
    ...                          # convert to Python integers and fractions, write one JSON line
\end{lstlisting}
\item \texttt{common.py} computes the terms and the five bounds from the records exactly as \texttt{rule} does:
\begin{lstlisting}[basicstyle=\ttfamily\scriptsize]
def terms(E, l, iii=1):
    out = []
    for p in E['primes']:
        q, k = p['q'], p['k']; e = w = t = 0
        if q != l:
            r = p['rf'] / (1 - Fraction(1, q * q)) if k >= 5 else p['rf']
            if abs(k) != 3 or iii: e = vl(r, l)
        elif l == 3: w = W3(k)
        if p['vj'] < 0: t = vl(-p['vj'], l)
        out.append((q, k, e, w, t))
    return out
def bounds(E, l):
    T1 = terms(E, l, 1); T0 = terms(E, l, 0)
    nc = sum(e + w + t for _, _, e, w, t in T1)
    h = E['ns3'] if l == 3 else 0
    main = sum(vl(phig(p['k']), l) for p in E['primes'])
    tw = 0
    for p in E['primes']:
        q, k = p['q'], p['k']
        if k >= 5 or (k <= -5 and q != 2): tw += vl(nidx(k), l)
        elif l == 3 and (abs(k) == 4 or (abs(k) == 2 and q != 2)): tw += 1
    thm = sum(e + t for q, _, e, w, t in T1 if q != l)
    return dict(rule=nc - h, nc=nc, iii=sum(e + w + t for _, _, e, w, t in T0),
                main=main, tw=tw, thm=thm)
\end{lstlisting}
Here \texttt{vl} is the $\ell$-adic valuation of a fraction, \texttt{W3} the weight of \Cref{conv:computation:code-idioms}, \texttt{phig} the order of
$\Phig_q$ read off the code and \texttt{nidx} the index $n$ of $\kod{I}_n$ or $\kod{I}_n^*$. The population filter compares the prime part of each \texttt{galrep} label
exactly, and CM is decided by $j$. The tables of \Cref{sec:computation:repro} are loops over the records with these functions.
\item \texttt{calib.py} computes \Cref{ex:computation:calibration-small,ex:computation:sharpness}; its core:
\begin{lstlisting}[basicstyle=\ttfamily\scriptsize]
def calib(P, l):
    data = [(bounds(E, l)['rule'], v(E, l)) for E in P]
    zero = [vv for b, vv in data if b == 0]; F = Counter(zero); n0 = len(zero)
    pos = [(b, vv) for b, vv in data if b > 0]
    ps = [sum(F[x] for x in F if x < b) / n0 for b, vv in pos]     # Fhat(v < b)
    Ehat = sum(ps); sd = math.sqrt(sum(p * (1 - p) for p in ps))
    Eeq = sum(F[b] / n0 for b, vv in pos)                            # expected equalities
    return Ehat, sd, Eeq, sum(vv < b for b, vv in pos), sum(vv == b for b, vv in pos)
\end{lstlisting}
The stratified version applies \texttt{calib} in each stratum; the permutation version shuffles $v$ within strata with \texttt{random.seed(0)}.
\item \texttt{tam.py} computes \Cref{ex:computation:repro-tam}, calling \texttt{ellinit(E, O(q\^{}30)).tate} at each qualifying prime (through the twist
by $q^*$ at a prime of type $\kod{I}_n^*$) and testing $\bar u^{(q-1)/\ell^k}=1$.
\item \texttt{moddeg\_check.py} compares \texttt{ellmoddegree} with \texttt{degphi} line by line (\Cref{ex:computation:moddeg-check}), and
\texttt{curvedata\_check.py} runs the cross-checks of \Cref{ex:computation:curvedata}.
\end{enumerate}
\end{remark}

\begin{convention}[a reproducibility checklist]\label{conv:computation:checklist}
Applied to the book's runs:
\begin{enumerate}[label=(\arabic*)]
\item \emph{Pin the data.} The commit of \texttt{ecdata} (first line) and the SHA-256 sums of the files read:
\begin{lstlisting}[basicstyle=\ttfamily\scriptsize]
commit 25cec5ecfec8b9f016eb1631ac633194c2bed39f  (4 February 2026)
b078b5226d276fbc64289833abfd8430e3f5969c657bc32f94db546e8f96e8fd  allbsd.00000-09999
259f3846329395b371e8079c77a6f1097adaebc98a054974573e241416efa968  allcurves.00000-09999
9e033f1d092509f280d23e2f2666577c96af9593911568ca9e9678c2e35b1d82  allisog.00000-09999
bc06d9bd10dbba36ddb84ad2fc407d2a71c8c89f22361befae8eda08aa18a377  alllabels.00000-09999
65d96cbab0de966409810b666e55b731515728c4d19904ea95b3750d48f6dd68  degphi.00000-09999
7faf8a6a92d85eba62dda9d025474a543bcd3ae61b8c489f581124f2ec34742d  galrep.00000-09999
89f4c055d668d7e7e55936bd5d391aaaa60df3133be8effc03394abcbc6ab7d5  alldegphi.00000-09999
b42836e901ac75874d13207faef0bdc47a3799571937b1099df56c57897e8040  curvedata.00000-09999
f1b7d7c310eb897ed4aad4da3e809b754e95be837cc4c1b404e1ed104b913223  classdata.00000-09999
11e4ebd69dbef2c6d9fce35c989d374f64a179e9d6ec13b5916cf453acd9298a  opt_man.00000-09999
\end{lstlisting}
The first six are the files of the counts; they were obtained independently of the repository and agree byte for byte with the commit.
\item \emph{Record the software and the settings.} PARI/GP 2.17.2, \texttt{cypari2} 2.2.4, Python 3.13; \texttt{parisize} $6.4\cdot10^7$ and \texttt{parisizemax} $2\cdot10^9$
bytes unless stated; random seed $0$ for the permutations.
\item \emph{Separate data, code and results}, and write results as plain files with the counts (the JSON records), not only as summaries.
\item \emph{Rerun from a clean copy} and compare the counts.
\item \emph{Cross-check independently.} Here: the GP loop against the Python bounds; \texttt{ellmoddegree} against \texttt{degphi}; \texttt{ns3} against
\texttt{galrep}; PARI's Euler factors against Lemma lem:euler; \texttt{ellglobalred} against \texttt{curvedata}; optimality flags against \texttt{degphi}.
\item \emph{State what was not checked.} Here: everything above conductor $10{,}000$; the paper's out-of-table curves, which are not published in the
paper; PARI 2.15; the contents of the paper's repository.
\end{enumerate}
\end{convention}

\begin{paperlink}\label{pl:computation:reproducibility}
The AI statement: ``The scripts, the result files and the pinned data commit are in the repository named in the appendix, so that every count
can be rerun''. Below conductor $10{,}000$ the pinned commit is the one used here, and every count of the paper that falls in this range has been
rerun from it: the Tamagawa counts of Section sec:tam (\Cref{ex:computation:repro-tam}), the two readings of the $\Ns$ violations
(\Cref{ex:computation:repro-ns}), the odd-part statistic (\Cref{ex:computation:odd-part}), the twist comparisons after Theorem thm:twist in their one-prime form
(\Cref{ex:computation:repro-twist}) and the data quoted in Remark rem:wild and Section sec:boundary, two of them with a different population than stated
(\Cref{sec:computation:populations,ex:computation:wild-two}).
\end{paperlink}
\section{Exercises}\label{sec:computation:exercises}

The exercises go from routine to harder. Each numerical answer was computed with PARI/GP 2.17.2 and the files of \Cref{conv:computation:checklist}.

\begin{exercise}\label{exer:computation:globalred}
Run \texttt{ellglobalred} on \curve{20a1} $=[0,1,0,4,4]$ and on \curve{1232l1} $=[0,0,0,32,-16]$. Read off the conductor, the bad primes, the Kodaira codes and
types, the conductor exponents and the $c_p$, and say which entries are the paper's \texttt{R[4][,1]} and \texttt{K[i][2]}. Why does the code $1$ never
occur in \texttt{R[5]}? (Answer: \curve{20a1}: $N=20$, $[[2,-4,0,3],[1,6,0,2]]$, so $\kod{IV}^*$ at $2$ with $f_2=2$, $c_2=3$ and $\kod{I}_2$ at $5$ with $c_5=2$;
\curve{1232l1}: $N=2^4\cdot7\cdot11$, $\kod{II}^*$ at $2$ with $f_2=4$, $\kod{I}_2$ at $7$, $\kod{I}_1$ at $11$, $\prod c_p=2$. The code $1$ is good reduction, and
\texttt{R[5]} lists only the bad primes.)
\end{exercise}

\begin{exercise}\label{exer:computation:sturm}
With \texttt{ellan}, check that $a_p(\curve{121a1})\equiv\chi_{-11}(p)a_p(\curve{11a1})\pmod3$ for every prime $p<500$, $p\ne3,11$ (use \texttt{kronecker(-11, p)}),
and that $a_n(\curve{121a1})\equiv a_n(\curve{121d1})\pmod3$ for all $n\le500$, with \curve{121d1} $=[0,-1,1,-40,-221]$. Explain why the second check, which is far
more than needed, proves the congruence for every $n$ (\texttt{mfsturm([121, 2])} returns $22=2\mu(121)/12$; \Cref{ch:modforms}; see \Cref{warn:computation:pitfalls}(h)), and why the first does not by
itself.
\end{exercise}

\begin{exercise}\label{exer:computation:euler}
Compute $P_q(X)$, $P_q(q^{-2})$ and $e_3(E,q)$ with \texttt{lfunsympow} and \texttt{lfuneuler} for \curve{147b1} $=[0,1,1,-114,473]$ at $7$, \curve{539d1} at
$7$ and \curve{726a1} at $11$, and compare with \texttt{euler\_by\_type}. (Answer: $1-7X$, $6/7$, $1$ (type $\kod{IV}^*$, $\chi_{-3}(7)=1$); $1+3X-21X^2-343X^3$, $360/343$, $2$;
$1+11X$, $12/11$, $1$.)
\end{exercise}

\begin{exercise}\label{exer:computation:allbounds}
Write a GP function \texttt{allbounds(E, l)} that returns, as a vector and in this order, $v_\ell(\degphi)$ and the bounds $B^{\mathrm{rule}}_\ell$,
$B^{\mathrm{nc}}_\ell$, $B^{\mathrm{III}}_\ell$, $B^{\mathrm{main}}_\ell$, $B^{\mathrm{tw}}_\ell$ of an optimal curve, reading $\#\Phig_q$ and $w_\ell(E,q)$ off the codes. Check that it returns $[1,1,1,1,0,1]$ for \curve{121a1} and $[2,2,3,2,2,2]$ for \curve{726a1} at $\ell=3$. Then show from the term-by-term comparison in the
proof of Proposition prop:forms that $B^{\mathrm{main}}_\ell\le B^{\mathrm{tw}}_\ell\le B^{\mathrm{nc}}_\ell$ for every curve and every odd $\ell$, and confirm it on the $38{,}042$ optimal
curves below $10{,}000$ at $\ell=3,5,7,11,13$.
\end{exercise}

\begin{exercise}\label{exer:computation:nonoptimal}
Using the files \texttt{alldegphi} and \texttt{opt\_man} of \Cref{ex:computation:curvedata}, explain the outputs \texttt{[1, 0]} and \texttt{[-1, 0]} of \texttt{rule} on
\curve{11a2} and \curve{11a3} at $\ell=5$: what is the degree of the minimal parametrisation of each, what is its Manin constant, and what does
\texttt{ellmoddegree} return? Which of these runs concerns a statement of the paper? (Answer: degrees $5$ and $5$, Manin constants $1$ and $5$,
\texttt{ellmoddegree} $5$ and $5/25=1/5$; none, since $E[5]$ is reducible and the curves are not optimal.)
\end{exercise}

\begin{exercise}\label{exer:computation:ns3-reducible}
Show that \texttt{ns3(E)} $=0$ whenever $E[3]$ is reducible, whatever the factorisation pattern of $\psi_3$. Find, below conductor $500$, an optimal curve with
\texttt{ns3(E)} $=1$ and complex multiplication, and say why Conjecture conj:euler and Table tab:ns exclude it. (Hint: a Galois-stable subgroup of order $3$
gives a rational root of $\psi_3$, so a factor of degree $1$. Answer: \curve{121b1}, with $j=-32768$; \curve{256a1}, with $j=8000$.)
\end{exercise}

\begin{exercise}\label{exer:computation:tablerow}
Using \texttt{table\_small.gp}, compute the rows $\ell=5$ and $\ell=7$ for $N<2000$ and compare with the values in \Cref{ex:computation:repro-uniform}. Then modify
\texttt{tablerow} to count, at $\ell=3$, the curves that violate the bound without the correction, and list them for $N<2000$. (Answer: \texttt{[5622, 915,
0, 845]} and \texttt{[5735, 468, 0, 446]}; the eight curves \curve{338d1}, \curve{338e1}, \curve{726a1}, \curve{726f1}, \curve{1210a1}, \curve{1210d1}, \curve{1210h1},
\curve{1210k1}.)
\end{exercise}

\begin{exercise}\label{exer:computation:ns-check}
For the eight curves of \Cref{exer:computation:tablerow}, check that the shortfall is $1$ and that each has a prime $p\equiv-1\pmod3$ with $t_3(E,p)\ge1$. Which
primes are they? Ten optimal curves with image $\Ns$ have conductor below $2000$; how many of them violate the bound without the correction, and what
fraction is that? (Answer: $p=2$ for all eight, and also $p=5$ for the four of conductor $1210$; $8$ of $10$.)
\end{exercise}

\begin{exercise}\label{exer:computation:tate}
Compute the Tate parameters of \curve{805b1} $=[1,-1,1,-163,-758]$ at $7$ and of \curve{874a1} $=[1,-1,0,-19,-13]$ at $19$, decide whether their unit parts are
cubes, and check that the refined bound of Section sec:tam equals $v_3(\degphi)$ in both cases. (Answer: $\bar u=1$ and $\bar u=12$, both cubes; the
rule gives $1$, the extra term is $1$ in both, and $\degphi=144$ in both.)
\end{exercise}

\begin{exercise}[dependence]\label{exer:computation:dependence}
Let $A_1,\dots,A_n$ be events of probability $p$ and $X$ the number that occur, so $E[X]=np$ in all cases. (a) If the $A_i$ are independent, show
$P(X=0)=(1-p)^n\le e^{-np}$. (b) If $A_1=\dots=A_n$, show $P(X=0)=1-p$. (c) Conclude that the expected count of \Cref{prop:computation:null-model}(i) needs
no independence but that every bound on $P(X=0)$ does, and say which kind of dependence among elliptic curves (twists, congruent newforms) pushes the
truth towards (b).
\end{exercise}

\begin{exercise}\label{exer:computation:f0}
The share of $v_3(\degphi)=0$ among the bound-$0$ curves at $\ell=3$ is $0.630$ on $N<5000$, $0.566$ on $5000\le N<10000$ and $0.532$ on the paper's working
set. Explain the trend, and why the paper's expected count of violations ($605{,}000$, about $70\%$ of the bounded curves) is close to the book's
proportion ($70.3\%$) although $\hat F(v=0)$ is smaller in the paper.
(Hint: $\hat F(v<b)$ for $b\ge2$ includes $v=1$, whose share grows with $N$.)
\end{exercise}

\begin{exercise}[multiple comparisons]\label{exer:computation:multiple}
Suppose $m=200$ false candidate inequalities each survive a working-set check with probability $\pi=0.05$, and a hold-out check with probability $\pi'=0.05$,
independently. How many false survivors are expected after the working set, and after the hold-out? If the same hold-out is then reused to choose among
$m'=200$ new candidates, how many false survivors does the reused hold-out let through, and what does a fresh hold-out achieve? (Answer: $10$, then
$0.5$; a reused hold-out acts as a working set and lets through about $10$; a fresh one brings this back to about $0.5$.)
\end{exercise}

\begin{exercise}\label{exer:computation:twistcheck}
Write \texttt{twist\_lemma\_check.gp}: for every optimal curve below conductor $2000$ and every odd bad prime $q$, compare the code and the conductor exponent
of the twist by $q^*$ at $q$ with Lemma lem:twist(ii), and check Lemma lem:twist(i) at the other primes. Extend the check to the potentially good types at
$q=3$ and compare with \Cref{ex:computation:twists}.
\end{exercise}

\begin{exercise}\label{exer:computation:fresh-holdout}
The paper's correction term $h_3$ was never blind-tested. Using the paper's numbers ($68$ violations without the correction on the hold-out, $23\%$ of
the $\Ns$ curves there), estimate how many $\Ns$ curves a fresh range of the same size would contain, and how many failures of the uncorrected bound it
would produce. Is that enough for a convincing test of ``the correction is exactly one on exactly these curves''? What else would the test need to
contain? (Answer: about $300$ $\Ns$ curves and $60$ to $70$ failures; the test would also need curves with image $\Ns$ and $3\nmid N$, where Theorem thm:main
does not apply only because of the image.)
\end{exercise}

\begin{exercise}\label{exer:computation:eis}
Reproduce the list of the $14$ Eisenstein curves at $\ell=3$ below $10{,}000$ with deficit $D=2$ (\Cref{ex:computation:repro-eis}), and compute $D$ for
\curve{14a1} by hand from its types ($\kod{I}_6$ at $2$, $\kod{I}_3$ at $7$) and $\degphi=1$. Then use \texttt{elltors} to confirm the torsion orders, and check that
none of the $49$ optimal curves with image \texttt{3Cs.1.1} has a point of order $9$.
\end{exercise}

\begin{exercise}[the protocol at $\ell=5$]\label{exer:computation:toy5}
Repeat \Cref{ex:computation:toy-protocol} at $\ell=5$: on $N<5000$, read off the minimum of $v_5(\degphi)$ over the curves of the base population with $5\mid N$,
potentially good reduction at $5$ and no other contribution, by the type at $5$; freeze the result; test it once on $5000\le N<10000$; and report honestly
whether it survived. Compare with the paper's ``at $q=\ell\ge5$ no term is needed''. (Answer: minimum $1$ for $\kod{II}^*$, $\kod{III}^*$, $\kod{IV}^*$ ($99$, $156$, $212$ curves) and
$0$ for $\kod{II}$, $\kod{III}$, $\kod{IV}$, $\kod{I}_0^*$; the same on the hold-out ($144$, $172$, $248$ curves for the starred types); see \Cref{ex:computation:twist-at-ell}.)
\end{exercise}

\begin{exercise}\label{exer:computation:oddpart}
Compute the odd-part statistic of \Cref{ex:computation:odd-part} on $N<5000$ and on $5000\le N<10000$, and explain why the paper's $1.85\%$ on $N\le300000$ is
so much smaller than $17.1\%$. (Answer: $25.5\%$ and $10.4\%$; as $N$ grows, degrees grow and the part of $\degphi$ not explained by local terms, the
Selmer group and the Petersson norm, is less and less often trivial.)
\end{exercise}

\section{Notes and references}\label{sec:computation:notes}

\emph{PARI/GP.} The functions are documented in the user's guide \cite{PARIusers}; the descriptions of \texttt{ellglobalred}, \texttt{elllocalred},
\texttt{ellmoddegree}, \texttt{ellweilcurve}, \texttt{lfuneuler}, \texttt{readstr} and \texttt{strsplit} used here were read from the documentation strings that
\texttt{cypari2} 2.2.4 carries for PARI 2.17.2, not from the printed manual of 2.15, and we cite no section numbers. The paper cites PARI/GP 2.15 \cite{PARI};
the book's runs are 2.17.2 \cite{PARI217} through \texttt{cypari2} \cite{cypari2}. Dokchitser's algorithm for $L$-values is \cite{Dokchitser04}; the method of
\texttt{ellmoddegree} is Watkins's \cite{Watkins02}.

\emph{Cremona's tables.} The book \cite{Cremona97} describes the algorithms and contains the tables to conductor $1000$; the \texttt{ecdata} repository
\cite{ecdata} extends them to conductor $500000$, and its file \texttt{docs/\allowbreak file-format.txt} documents the formats of most files (not
\texttt{curvedata}, \texttt{classdata} and \texttt{opt\_man}, whose columns we read from the repository's scripts and from the data). The LMFDB \cite{LMFDB} uses
the same curves with its own labels. Sutherland's labels of mod $\ell$ images are defined in \cite{Sutherland16}, and Zywina's classification \cite{Zywina15} is
the source of the list of images at $\ell=3$ that Proposition prop:ns uses. Modular degrees and their computation: \cite{Zagier85,Watkins02} and
\cite[Chapter II]{Cremona97}; the congruence number of \curve{99A1} is in \cite[Table 1]{ARS12}. Tate's algorithm and the Kodaira types are in
\cite[Chapter IV]{SilvermanATAEC}.

\emph{Statistics.} The vocabulary is that of Appendix A (\Cref{ch:prereq}); for expectation, Chebyshev's inequality and the law of large numbers see
Feller \cite{Feller68} (the chapters on random variables and expectation and on the law of large numbers); for hypothesis tests, permutation tests and
multiple testing see Wasserman \cite{Wasserman04} (the chapters on hypothesis testing and $p$-values); Riley, Hobson and Bence \cite{RHB06} cover the
reader's background. Chapter titles are cited from memory and are pointers we have not checked against the printed editions.

\emph{What was verified, and how.} Every number marked as the book's was produced with PARI/GP 2.17.2 through \texttt{cypari2} 2.2.4 and Python 3.13 on
Linux, from the files of commit \texttt{25cec5e} of \texttt{ecdata} listed in \Cref{conv:computation:checklist}; the printed GP code was loaded with
\texttt{read} and run as printed, and every printed output was copied from the run. The paper's numbers for its large ranges are quoted from its text
and were not reproduced. Not checked: PARI 2.15; the paper's repository; the out-of-table curves; anything above conductor $10{,}000$.

\emph{Corrections to the paper recorded here.} Remark rem:wild: the counts $1{,}143$ and $2{,}283$ of cubic factors are over all curves, not optimal curves
without CM ($534$ and $1{,}308$), and its $58\%$ is a share among optimal curves ($61\%$ over all curves) (\Cref{ex:computation:wild-two}). Section sec:boundary:
$44$, not $28$, of the $90$ optimal $\Ns$ curves have $3\mid N$, and $13$ optimal CM curves with image $N_s(\F_3)$ are left out. Section sec:forms: the
$1{,}170$ failures of the hypothesis at $\ell=3$ are those without CM; $20$ CM curves below $10{,}000$ also fail ($10$ optimal)
(\Cref{sec:computation:populations}). Section sec:anomaly: the two counts $36$ and $24$ differ by the terms $t_3(E,3)$ at $q=3$, not by the weights $w_3$
(\Cref{ex:computation:repro-ns}). The text of Section sec:data leaves two questions of protocol open (\Cref{rem:computation:test-status}).
\end{sloppypar}
